\documentclass[11pt,a4paper]{article}
\usepackage[utf8]{inputenc}
\usepackage[T1]{fontenc}
\usepackage[english]{babel}
\usepackage{amsmath,amssymb,amsthm}
\usepackage{geometry}
\usepackage{titlesec}
\usepackage{fancyhdr}
\usepackage{microtype}

\geometry{margin=1in}

\pagestyle{fancy}
\fancyhf{}
\fancyhead[C]{\small Cayley Mathematical Papers, Vol.\ XIII, pp.\ 51--100}
\fancyfoot[C]{\thepage}
\renewcommand{\headrulewidth}{0.4pt}

\titleformat{\section}{\large\bfseries}{\thesection.}{0.5em}{}
\titleformat{\subsection}{\normalsize\bfseries}{\thesubsection.}{0.5em}{}

\theoremstyle{plain}
\newtheorem{theorem}{Theorem}

\DeclareMathOperator{\Disc}{Disct.}

\begin{document}

\begin{center}
\Large\textbf{The Collected Mathematical Papers of}\\[4pt]
\LARGE\textbf{Arthur Cayley}\\[8pt]
\Large\textbf{Volume XIII}\\[12pt]
\large\textbf{Pages 51--100}
\end{center}

\vspace{1em}
\hrule
\vspace{1.5em}

% ========================================================================
% PAPER 896 (continued from p. 31): A TRANSFORMATION IN ELLIPTIC FUNCTIONS
% ========================================================================

\section*{896. A Transformation in Elliptic Functions.}
\addcontentsline{toc}{section}{896. A Transformation in Elliptic Functions.}

\noindent\textit{[From the Messenger of Mathematics, vol.\ xix.\ (1890), pp.\ 30--33.]}

\vspace{0.5em}
\noindent\textbf{[Continuation from p.\ 31]}

which is
$$
\frac{1}{\sqrt{vw}},
$$
and similarly the whole coefficient of $dx_2^2$ is
$$
\frac{1}{\sqrt{vw}}.
$$
Hence the terms in $dx_1$, $dx_2$ are together
$$
(x_1 - x_2)\{D_x\sqrt{J(A)} + B_x\sqrt{J(E)}\}.
$$
and since the terms in $dy_1$, $dy_2$ are of the like form, we have
$$
\frac{dz}{\sqrt{(1 - Lz - J)}} = \frac{1}{\sqrt{J(A)}} \cdot \frac{x_1\,dx_2 - x_2\,dx_1}{x_1 - x_2} - \frac{1}{\sqrt{J(E)}} \cdot \frac{y_1\,dy_2 - y_2\,dy_1}{y_1 - y_2},
$$
and combining herewith the foregoing value
$$
\frac{dz}{\sqrt{(1 - Lz - J)}} = \frac{x_1\,dx_2 - x_2\,dx_1}{x_1 - x_2} - \frac{y_1\,dy_2 - y_2\,dy_1}{y_1 - y_2},
$$
we have the required formula
$$
\frac{dz}{\sqrt{(1 - Lz - J)}} = \frac{x_1\,dx_2 - x_2\,dx_1}{x_1 - x_2} - \frac{y_1\,dy_2 - y_2\,dy_1}{y_1 - y_2}.
$$

As a very simple verification, suppose $a = 1$, $b = c = d = e = 0$; then
$$
(A, B, C, D, E) = (x_1^4,\; x_1^3 y_1,\; x_1^2 y_1^2,\; x_1 y_1^3,\; y_1^4),
$$
and if $\lambda = x_1 y_2 - x_2 y_1$, as before, then
$$
z = \frac{x_1^2 y_1^2}{(x_1 y_2 - x_2 y_1)^2} = \theta^2,
$$
where
$$
\theta = \frac{-x_1 y_1}{x_1 y_2 - x_2 y_1}, \quad \text{or} \quad \frac{1}{\theta} = \frac{x_2}{x_1} - \frac{y_2}{y_1}.
$$
Also $I = 0$, $J = 0$, and consequently
$$
\sqrt{(1 - Lz - J)} = \sqrt{(1 - z^2)},
$$
which, in virtue of the foregoing values of $A$, $E$, is
$$
\frac{dz}{\sqrt{(1 - z^2)}} = \frac{x_1\,dx_2 - x_2\,dx_1}{x_1 - x_2} - \frac{y_1\,dy_2 - y_2\,dy_1}{y_1 - y_2}.
$$

I remark that an even more simple transformation from the general quartic radical to the Weierstrassian cubic radical was obtained by Hermite; this is alluded to in my ``Note sur les Covariants, \&c.,'' Crelle, t.\ L.\ (1855), pp.\ 285--287, [135], and is given with a demonstration in my paper ``Sur quelques formules pour la transformation des int\'egrales elliptiques,'' Crelle, t.\ LV.\ (1858), pp.\ 15--24, [235], see No.\ iv.; viz. from the identical relation $JU^3 - IU^2H + 4H^3 = -\Phi^2$, which connects the covariants of a quartic function, it at once follows that if
$$
U = (a, b, c, d, e \!\!\raisebox{-2pt}{$\bigm|$}\, x, 1)^4,
$$
and $H$ is the Hessian hereof,
$$
H = (ac - b^2,\; ad - bc,\; ae + 2bd - 3c^2,\; 2(be - cd),\; ce - d^2 \!\!\raisebox{-2pt}{$\bigm|$}\, x, 1)^4;
$$
then writing
$$
z = \frac{H}{U},
$$
we have
$$
\frac{dz}{\sqrt{(-4z^3 + Iz - J)}} = \frac{2\,dx}{\sqrt{(a, b, c, d, e \!\!\raisebox{-2pt}{$\bigm|$}\, x, 1)^4}},
$$
which is the transformation in question.

\newpage
% ========================================================================
% PAPER 897: SUR LES RACINES D'UNE EQUATION ALGEBRIQUE
% ========================================================================

\section*{897. Sur les Racines d'une \'Equation Alg\'ebrique.}
\addcontentsline{toc}{section}{897. Sur les Racines d'une \'Equation Alg\'ebrique.}

\noindent\textit{[From the Comptes Rendus de l'Acad\'emie des Sciences de Paris, t.\ cx.\ (Janvier--Juin, 1890), pp.\ 174--176, 215--218.]}

\vspace{0.5em}
\noindent\textbf{[p.\ 33]}

Soit $f(u)$ une fonction rationnelle et enti\`ere avec des coefficients r\'eels ou imaginaires, de l'ordre $n$; en supposant que l'\'equation $f'(u) = 0$, de l'ordre $n - 1$, ait $n - 1$ racines, je d\'emontre que l'\'equation $f(u) = 0$ aura $n$ racines. Pour cela, soit
$$
f(u) = f(x + iy) = P + iQ;
$$
je suppose que $c$ d\'enote une quantit\'e positive donn\'ee, et je consid\`ere la surface $c - z = P^2 + Q^2$, en attribuant \`a la coordonn\'ee $z$ des valeurs positives; c'est seulement pour avoir des maxima au lieu de minima, et pour faciliter ainsi l'exposition, que je prends cette surface au lieu de $z = P^2 + Q^2$. On peut donner \`a $c$ une valeur si grande que la courbe $c = P^2 + Q^2$ soit une courbe ferm\'ee qui ne se coupe pas, c'est-\`a-dire un contour simple: cela \'etant, on peut se figurer ce contour comme la ligne de rivage d'une \^ile montagneuse; la valeur de $z$ est au plus $= c$, et, en donnant \`a $z$ une valeur plus petite, $= b$, on a le contour qui correspond \`a l'altitude $b$: \'evidemment, les contours qui correspondent \`a des altitudes diff\'erentes ne se coupent pas. Il s'agit de prouver que l'\^ile a pr\'ecis\'ement $n$ sommets, chacun de l'altitude $c$.

J'\'ecris
$$
X = \frac{dP}{dx} = \frac{dQ}{dy}, \quad Y = \frac{dP}{dy} = -\frac{dQ}{dx};
$$
donc
$$
\frac{dQ}{dx} = -Y, \quad \frac{dQ}{dy} = X,
$$
et, de plus,
$$
\frac{dX}{dx} = \frac{d^2 P}{dx^2}, \quad \frac{dX}{dy} = \frac{d^2 P}{dx\,dy}, \quad \frac{dY}{dx} = \frac{d^2 P}{dx\,dy}, \quad \frac{dY}{dy} = \frac{d^2 P}{dy^2}.
$$

\vspace{0.5em}
\noindent\textbf{[p.\ 34]}

et de l\`a
$$
\frac{d^2 z}{dx^2} = -2(X^2 + Y^2) - 2P\frac{d^2 P}{dx^2} - 2Q\frac{d^2 Q}{dx^2}, \quad \text{etc.}
$$

Cela \'etant, nous avons
$$
\frac{dz}{dx} = -2\left(P\frac{dP}{dx} + Q\frac{dQ}{dx}\right), \quad \frac{dz}{dy} = -2\left(P\frac{dP}{dy} + Q\frac{dQ}{dy}\right);
$$
on aura un plan tangent horizontal pour $P = 0$, $Q = 0$, ou pour $X = 0$, $Y = 0$; les valeurs $P = 0$, $Q = 0$, appartiennent \`a la valeur $c$ de $z$ et correspondent \`a des sommets de montagne de cette altitude $c$; les valeurs $X = 0$, $Y = 0$ correspondent \`a des sommets de col. En effet, nous avons
$$
\frac{d^2 z}{dx^2} = -2(X^2 + Y^2) - 2P\frac{dX}{dx} - 2Q\frac{dY}{dx}, \quad \text{etc.},
$$
done
$$
\frac{d^2 z}{dx^2}\frac{d^2 z}{dy^2} - \left(\frac{d^2 z}{dx\,dy}\right)^2 = (X^2 + Y^2)^2 - (P^2 + Q^2)(a^2 + A^2),
$$
valeur positive pour $P = 0$, $Q = 0$; n\'egative pour $X = 0$, $Y = 0$.

\`A pr\'esent, nous avons $f'(x + iy) = X - iY$; donc, en supposant que l'\'equation $f'(x + iy) = 0$ ait $n - 1$ racines, il y aura $n - 1$ syst\`emes de valeurs r\'eelles de $x$, $y$ qui satisfont aux \'equations $X = 0$, $Y = 0$; pour chaque syst\`eme, il y aura des valeurs d\'etermin\'ees de $P$, $Q$ et de l\`a aussi de $z$; ces valeurs de $z$ seront en g\'en\'eral diff\'erentes. Ainsi il y aura dans l'\^ile $n - 1$ cols, dont les altitudes seront en g\'en\'eral diff\'erentes: soient $c_1$, $c_2$, \ldots, $c_{n-1}$ ces altitudes, commen\c{c}ant avec la plus petite.

Pour $b = 0$, nous avons un contour simple, et de m\^eme pour une valeur quelconque plus petite que $c_1$; mais, pour $z = c_1$, nous avons un col; le contour est une courbe, figure de 8; en donnant \`a $z$ une valeur un peu plus grande, le contour se divise en deux courbes ferm\'ees, ou bien contours simples, ext\'erieurs l'un \`a l'autre. Pour $z = c_2$, nous avons encore un col; l'un des contours simples s'est chang\'e en figure de 8, et, pour une valeur un peu plus grande, le contour se divise en trois contours simples, chacun ext\'erieur aux autres. En continuant de cette mani\`ere, on a, pour $c_{n-1}$, le col le plus haut; le contour est compos\'e de $n - 2$ contours simples et d'une figure de 8; et, en donnant \`a $z$ une valeur un peu plus grande, on obtient un contour compos\'e de $n$ contours simples, chacun ext\'erieur aux autres. Enfin, en

\vspace{0.5em}
\noindent\textbf{[p.\ 35]}

faisant cro\^itre $z$, chacun des contours simples doit se r\'eduire \`a un point, c'est-\`a-dire qu'il doit y avoir pr\'ecis\'ement $n$ sommets de montagne; mais il n'y a pas de sommet de montagne, sinon pour la valeur $z = c$; donc il y a pr\'ecis\'ement $n$ sommets de montagne, chacun de l'altitude $c$. On suppose toujours que l'\'equation $f'(x + iy) = 0$ n'ait pas de racines \'egales, mais il peut bien arriver que deux ou plusieurs des valeurs $c_1$, $c_2$, \ldots, $c_n$ deviennent \'egales; la d\'emonstration est tr\`es peu chang\'ee, en donnant \`a $z$ une valeur un peu plus grande que celle qui correspond \`a l'altitude des cols d'altitude \'egale; le contour se divise toujours en contours simples, ext\'erieurs chacun aux autres.

Il va sans dire que cette d\'emonstration repose sur les m\^emes principes que celles de Gauss et de Cauchy.

Je reprends la th\'eorie des racines de l'\'equation $f(u) = 0$; au lieu de la surface $c - z = P^2 + Q^2$, il convient de consid\'erer la surface $(c - z)^2 = P^2 + Q^2$, en faisant attention seulement aux valeurs de $z$ positives et pas plus grandes que $c$. La th\'eorie est tr\`es peu chang\'ee; les contours sont les m\^emes qu'auparavant, mais ils appartiennent \`a des altitudes diff\'erentes; et, au lieu de maxima $z = c$ pour $P = 0$, $Q = 0$, on a des points coniques, c'est-\`a-dire, dans l'\^ile montagneuse, au lieu d'un sommet arrondi de montagne, on a un c\^one ou un pic.

Mais avec la nouvelle surface, on construit graphiquement l'approximation de Newton: partant d'une valeur r\'eelle ou imaginaire approximative $u$, on obtient la nouvelle valeur
$$
u_1 = u - \frac{f(u)}{f'(u)}.
$$
Je repr\'esente $u$ par le point $(x, y, z)$ de la surface $(c - z)^2 = P^2 + Q^2$, ou le point $(x, y)$ du plan des sommets $z = c$; et de m\^eme, $u_1$ par le point $(x_1, y_1, z_1)$ de la surface, ou $(x_1, y_1)$ du plan des sommets: cela \'etant, si, par le point $(x, y, z)$ de la surface, on m\`ene la droite de plus grande pente (droite tangente \`a la surface et perpendiculaire au contour), cette droite rencontrera le plan des sommets en un point $(x_1, y_1)$, et l'on obtient ainsi le point $(x_1, y_1, z_1)$ de la surface, qui repr\'esente la valeur cherch\'ee $u_1$. En particulier, si les coefficients de $f(u)$ sont r\'eels, on a $Q = 0$; l'\'equation $(c - z)^2 = P^2 + Q^2$ devient
$$
(c - z)^2 = P^2,
$$
c'est-\`a-dire
$$
c - z = \pm P,
$$
ou enfin
$$
c - z = \pm f(x),
$$
et la section verticale de l'\^ile est form\'ee par des parties de ces deux courbes sym\'etriques: pour la th\'eorie g\'eom\'etrique, on peut \'evidemment y substituer la seule courbe $c - z = f(x)$.

\vspace{0.5em}
\noindent\textbf{[p.\ 36]}

J'ai propos\'e, il y a plus de dix ans (Amer.\ Math.\ Journ., t.\ II., 1879, [694]), le probl\`eme que je nomme ``The Newton-Fourier imaginary Problem,'' et dans une Note (Quart.\ Math.\ Journ., t.\ xvi., 1879, [736]), ``Application of the Newton-Fourier method to an imaginary root of an equation,'' j'ai consid\'er\'e le cas d'une \'equation quadratique. Pour l'\'equation $u^2 - 1 = 0$, on a
$$
u_1 = \frac{1}{2}\left(u + \frac{1}{u}\right), \quad \text{or} \quad x_1 + iy_1 = \frac{1}{2}\left(x + iy + \frac{1}{x + iy}\right);
$$
cela donne
$$
x_1 = \frac{1}{2}\left(x + \frac{x}{x^2 + y^2}\right), \quad y_1 = \frac{1}{2}\left(y - \frac{y}{x^2 + y^2}\right),
$$
et de l\`a
$$
\frac{u_1 - 1}{u_1 + 1} = \left(\frac{u - 1}{u + 1}\right)^2.
$$
Cette derni\`ere \'equation a rapport aux deux racines $+1$ et $-1$, et, quoiqu'elle donne les r\'esultats les plus \'el\'egants, cependant, en vue de la th\'eorie g\'en\'erale, il vaut mieux consid\'erer l'\'equation $u_1 - 1 = \frac{1}{2}(u - 1)^2$ qui se rapporte \`a la seule racine $+1$.

Je remarque d'abord que la formule originale $u_1 = \frac{1}{2}(u + \frac{1}{u})$ donne
$$
\frac{x_1}{y_1} = \frac{x}{y} \cdot \frac{x^2 + y^2 + 1}{x^2 + y^2 - 1},
$$
donc les valeurs de $x$ et $x_1$ seront \`a la fois positives ou n\'egatives, et ainsi l'on peut ne faire attention qu'aux valeurs positives. Cela \'etant, nous avons
$$
x_1 - 1 + i y_1 - 1 = \frac{(x + iy - 1)^2}{2(x + iy)}.
$$
D\'esignons par $A$ le point $(0, 1)$, par $B$ le point $(0, -1)$, par $O$ le point $(0, 0)$; et aussi par $P$ le point $(x, y)$, et de m\^eme

\vspace{0.5em}
\noindent\textbf{[p.\ 37]}

par $P_1$ le point $(x_1, y_1)$. En consid\'erant l'expression de $u_1 - 1$, on voit que la distance du point $P_1$ au point $A$ est le carr\'e de la distance du point $P$ au point $A$, divis\'e par le double de la distance du point $P$ au point $O$. Cela \'etant, soit $\rho$ la distance $AP$, et $\rho_1$ la distance $AP_1$: on a
$$
\rho_1 = \frac{\rho^2}{2OP}, \quad \text{or} \quad AP_1 = \frac{(AP)^2}{2OP}.
$$
Ainsi, pour le point $P$ situ\'e sur l'axe des $x$, on a $AP_1 < AP$; c'est-\`a-dire que l'approximation est r\'eguli\`ere: et l'on peut m\^eme (dans un sens qui sera expliqu\'e plus bas, mais qui n'est pas le sens le plus naturel) dire que l'approximation est r\'eguli\`ere. En effet, on n'a pas toujours $AP_1 < AP$, et ainsi, dans le sens le plus naturel, l'approximation n'est pas toujours r\'eguli\`ere. Pour \'etudier cela, j'\'ecris $AP_1 = AP$; cela donne $AP = 2OP$, ou, ce qui est la m\^eme chose, $x^2 + y^2 + \frac{1}{4} = \frac{1}{4}$, c'est-\`a-dire que le point $P$ sera situ\'e sur le cercle, centre $x = -\frac{1}{2}$ et rayon $= \frac{1}{2}$; en ne faisant attention qu'aux valeurs positives de $x$, on a un segment compris entre l'axe des $y$ et un arc par les points $(x = 0,\; y = \pm \frac{1}{2})$, $(x = \frac{1}{2},\; y = 0)$. Si le point $P$ est sur l'arc, on aura $AP_1 = AP$; si $P$ est en dedans du segment, alors $AP_1 > AP$; si $P$ est en dehors du segment, $AP_1 < AP$.

Mais, en supposant $P$ en dehors du segment, et ainsi $AP_1 < AP$, il peut bien arriver que $P_1$ soit en dedans du segment, et, cela \'etant, on aura $AP_2 > AP_1$, et l'approximation ne sera pas r\'eguli\`ere. Mais, en consid\'erant le cercle $x^2 + y^2 - \frac{1}{2}x = \frac{1}{4}$, lequel est le cercle, centre $A$ et rayon $\frac{1}{2}$, qui touche le segment au point $(x = \frac{1}{2},\; y = 0)$, alors, en supposant que le point $P$ soit en dedans de ce cercle, on aura $AP_1 < AP$, le point $P_1$ sera aussi en dedans du cercle, et ainsi en dehors du segment; et les points successifs $P$, $P_1$, $P_2$, \ldots\ approcheront continuellement du point $A$; l'approximation sera dans ce cas r\'eguli\`ere.

Il y a ainsi trois r\'egions: le segment, le cercle $x^2 + y^2 - \frac{1}{2}x = \frac{1}{4}$ et le r\'esidu du demi-plan; on pourrait les nommer r\'egions noire, blanche et grise respectivement. C'est seulement pour un point $P$ \`a l'int\'erieur de la r\'egion blanche que l'approximation est certainement r\'eguli\`ere.

Nous venons de consid\'erer en effet les cercles qui ont pour centre le point $A$; $AP_1 < AP$ veut dire que le point $P$ est situ\'e sur un cercle plus grand, et $P_1$ sur un cercle plus petit; mais, au lieu de ces cercles concentriques, consid\'erons des cercles quelconques qui entourent le point $A$, sans se couper les uns les autres; et convenons de dire que c'est un bon pas quand on passe du point $P$ sur un cercle plus grand \`a un point $P_1$ sur un cercle plus petit: avec cette convention on aura, en g\'en\'eral, trois r\'egions, lesquelles cependant ne seront pas les m\^emes comme auparavant. En particulier, si nous consid\'erons les cercles $AP = k \cdot BP$ ($k$ une constante quelconque plus petite que l'unit\'e), alors il n'y a pas de r\'egion noire, ou, si l'on veut, la r\'egion noire se r\'eduit \`a la seule droite $y = 0$; donc il n'y a pas non plus de r\'egion grise, et le demi-plan entier est r\'egion blanche, c'est-\`a-dire, dans le sens que je viens d'expliquer, l'approximation est toujours r\'eguli\`ere. En effet, c'est l\`a la th\'eorie \`a laquelle on est conduit au moyen de l'\'equation $\frac{u_1 - 1}{u_1 + 1} = \left(\frac{u - 1}{u + 1}\right)^2$ ci-dessus mentionn\'ee.

En parlant de cercles, j'ai fait une restriction qui n'est nullement n\'ecessaire; j'aurais pu parler d'ovales, de forme quelconque, qui entourent le point $A$ sans se couper les uns les autres.

J'esp\`ere appliquer cette th\'eorie au cas d'une \'equation cubique, mais les calculs sont beaucoup plus difficiles.

\newpage
% ========================================================================
% PAPER 898: SUR L'EQUATION MODULAIRE POUR LA TRANSFORMATION DE L'ORDRE 11
% ========================================================================

\section*{898. Sur l'\'Equation Modulaire pour la Transformation de l'Ordre 11.}
\addcontentsline{toc}{section}{898. Sur l'\'Equation Modulaire pour la Transformation de l'Ordre 11.}

\noindent\textit{[From the Comptes Rendus de l'Acad\'emie des Sciences de Paris, t.\ cxi.\ (Juillet--D\'ecembre, 1890), pp.\ 447--449.]}

\vspace{0.5em}
\noindent\textbf{[p.\ 38]}

L'\'equation en $u$, $v$, en y \'ecrivant $u = x$, $v = y$, est
\begin{multline*}
x^{12} + y^{12} - x^5 y^5 + 11x^6 y^6 (x^5 + y^5) - 66x^4 y^4 (x^5 + y^5)\\
+ 11x^3 y^3 (x^5 + y^5) - 11x^2 y^2 (x^5 + y^5) - 11xy(x^5 + y^5) - 1 = 0.
\end{multline*}

Selon un r\'esultat trouv\'e par H.\ J.\ S.\ Smith, pour la transformation de l'ordre $p$, la courbe est de l'ordre $2p$, et il y a \`a l'origine un point double, \`a l'infini deux points singuliers \'equivalents chacun \`a $\frac{1}{2}(p - 1)(p - 2)$ points doubles, et de plus $(p - 1)(p - 3)$ points doubles. Au cas $p = 11$, le nombre de ces derniers points doubles est donc $= 80$. Cela s'accorde avec l'expression
$$
D = x^{116}(1 - x^8)^{10}(16x^{16} - 31x^8 + 16)^8(x^{64} - 301960x^{56} + \ldots + 1)^2,
$$

\vspace{0.5em}
\noindent\textbf{[p.\ 39]}

trouv\'ee par M.\ Hermite pour le discriminant de la fonction; et l'on voit ainsi que les valeurs de $x$, qui correspondent aux quatre-vingts points doubles, sont donn\'ees par les \'equations
$$
16x^{16} - 31x^8 + 16 = 0,
$$
$$
x^{64} - 301960x^{56} + \ldots + 1 = 0.
$$
Je ne consid\`ere que les seize points doubles donn\'es par la premi\`ere \'equation. Cette \'equation donne
$$
x^8 = \tfrac{1}{32}(31 \pm 3\sqrt{7}), \quad x^4 = \tfrac{1}{4}(3 \pm \sqrt{7}),
$$
il y a ainsi quatre points doubles, pour lesquels les valeurs de $x$ sont
$$
x = \pm\sqrt{\frac{3 + \sqrt{7}}{4}}, \quad x = \pm i\sqrt{\frac{3 - \sqrt{7}}{4}}.
$$
Je trouve que les valeurs correspondantes de $y$ sont $y = -mx$, savoir que les quatre points sont situ\'es sur la droite $y = -\frac{1+i}{2}x$, et, en changeant successivement les signes de $i$ et $\sqrt{2}$, on voit ainsi que les seize points sont situ\'es, quatre \`a quatre, sur les droites
$$
y = -\frac{1+i}{2}x, \quad y = \frac{i-1}{2}x, \quad y = \frac{1+i}{2}x, \quad y = \frac{1-i}{2}x.
$$
J'\'ecris, pour abr\'eger,
$$
m = \frac{1+i}{2} \quad \text{(donc } m^4 = -1 \text{)},
$$
done
$$
y = -mx.
$$
En \'ecrivant $y = mx$ dans l\'equation et en rejetant le facteur $x^2$, puis en \'ecrivant $\rho = mp$, l\'equation se pr\'esente sous la forme
$$\small
\left.\begin{array}{r@{\,}l@{\quad}r}
m^{10}\rho^{10} & .\quad 32m^{11} \\
{}+\; m^{8}\rho^{8} & .\quad 88m^{9} \\
{}+\; m^{7}\rho^{7} & .\quad 44m^{10} - 44m^{6} \\
{}+\; m^{6}\rho^{6} & .\quad -22m^{11} + 132m^{7} - 22m^{3} \\
{}+\; m^{5}\rho^{5} & .\quad m^{12} + 165m^{8} - 165m^{4} - 1 \\
{}+\; m^{4}\rho^{4} & .\quad 22m^{9} - 132m^{5} + 22m \\
{}+\; m^{3}\rho^{3} & .\quad 44m^{6} - 44m^{2} \\
{}+\; m^{2}\rho^{2} & .\quad -88m^{3} \\
{}+\; 1 & .\quad -32m
\end{array}\right\} = 0,
$$
ou les coefficients ne contiennent que les puissances $m^{21}$, $m^{17}$, $m^{13}$, $m^9$, $m^5$, $m^1$ de $m$ et se r\'eduisent ainsi \`a des multiples de $m$; il y a aussi un facteur num\'erique 8, et, en divisant par $-8m$, l\'equation devient
$$4\rho^{10} - 11\rho^8 - 11\rho^7 + 22\rho^6 + 41\rho^5 + 22\rho^4 - 11\rho^3 - 11\rho^2 + 4 = 0;$$
cette \'equation est de la forme
$$
(2\rho^2 + 3\rho + 2)^2(\rho^6 - 3\rho^5 + 2\rho^4 + \rho^3 + 2\rho^2 - 3\rho + 1) = 0.
$$
La droite $y = mx$ a donc, avec la courbe, quatre intersections doubles
$$
\rho = \frac{-3 \pm \sqrt{7}}{4}, \quad \rho^6 - 3\rho^5 + 2\rho^4 + \rho^3 + 2\rho^2 - 3\rho + 1 = 0,
$$
c'est-\`a-dire
$$
2x^2 = \frac{-3 \pm \sqrt{7}}{4m^2}, \quad \text{ou} \quad x^2 = \frac{-3 \pm \sqrt{7}}{4\sqrt{2}}i;
$$
on d\'emontre sans peine que la droite n'est pas une tangente, et ces valeurs correspondent ainsi \`a des points doubles de la courbe, c'est-\`a-dire qu'il y a sur la droite $y = \frac{1+i}{2}x$ quatre points doubles. R\'eciproquement, cette valeur de $x^2$ conduit au facteur $(16x^{16} - 31x^8 + 16)^2$ du d\'eterminant de l'\'equation modulaire.

\newpage
% ========================================================================
% PAPER 899: SUR LES SURFACES MINIMA
% ========================================================================

\section*{899. Sur les Surfaces Minima.}
\addcontentsline{toc}{section}{899. Sur les Surfaces Minima.}

\noindent\textit{[From the Comptes Rendus de l'Acad\'emie des Sciences de Paris, t.\ cxi.\ (Juillet--D\'ecembre, 1890), pp.\ 953, 954.]}

\vspace{0.5em}
\noindent\textbf{[p.\ 41]}

On peut g\'en\'eraliser tant la d\'efinition que la construction de ces surfaces, en substituant pour le cercle imaginaire \`a l'infini une conique ou m\^eme une surface quadrique quelconque.

Je rappelle que, dans la th\'eorie ordinaire, une surface minima est une surface telle qu'un point quelconque de la surface est situ\'e \`a mi-chemin entre les deux centres de courbure, et qu'une telle surface est le lieu des points \`a mi-chemin entre les deux points situ\'es respectivement sur deux courbes de longueur nulle quelconques.

Or on peut rattacher la notion d'une courbe de longueur nulle \`a celle d'une courbe de poursuite. Pour expliquer cela, j'observe que, dans le plan, en supposant comme \`a l'ordinaire que le li\`evre coure selon une droite et que le chien et le li\`evre courent avec des vitesses uniformes, la courbe de poursuite est une courbe d\'etermin\'ee; mais si les vitesses varient arbitrairement, alors la d\'efinition exprime seulement que la courbe est une courbe plane. Mais, dans l'espace, si au lieu d'une droite on consid\`ere une courbe plane ou \`a double courbure (disons une directrice) quelconque, alors, quelles que soient les vitesses, la d\'efinition pr\'ecise toujours la courbe, savoir: on a toujours pour courbe de poursuite une courbe dont chaque tangente rencontre la courbe directrice. Au lieu d'une courbe, on peut avoir une surface directrice; dans ce cas, le nom est moins applicable, n\'eanmoins je le retiens, et je dis que la courbe de poursuite est une courbe dont chaque tangente touche la surface directrice. Nous avons, de cette mani\`ere, la d\'efinition d'une courbe de poursuite par rapport \`a une courbe ou surface directrice quelconque.

\`A pr\'esent, au lieu du cercle imaginaire \`a l'infini, consid\'erons une conique quelconque, l'absolue: on \'etablit, comme on sait, par rapport \`a cette conique, la notion de la perpendicularit\'e, et ainsi les notions d'une normale et des centres de courbure

\vspace{0.5em}
\noindent\textbf{[p.\ 42]}

ne cessent pas de subsister. On peut donc consid\'erer une surface telle que chaque point de la surface soit l'harmonicale par rapport aux deux centres de courbure du point de rencontre de la normale avec le plan de l'absolue: on a ainsi ce que je nomme une surface quasi-minima. Il va sans dire qu'il faut modifier convenablement la notion m\'etrique d'une aire minima pour qu'elle soit applicable \`a cette nouvelle surface.

Pour construire la surface, on prend par rapport \`a l'absolue deux courbes de poursuite quelconques, et puis sur la droite, men\'ee par deux points quelconques de ces courbes respectivement, l'harmonicale par rapport \`a ces deux points du point de rencontre de la droite avec le plan de l'absolue: le lieu de ce point harmonical sera la surface quasi-minima.

Il para\^it permis de substituer pour la conique absolue une surface quadrique quelconque, que je nomme aussi l'absolue: on a, comme on sait, les notions de la normale et des centres de courbure. Pour la surface quasi-minima, le point sur la surface sera l'un des points doubles (foyers) de l'involution form\'ee par les deux centres de courbure et les deux points de rencontre de la normale avec l'absolue; et de m\^eme pour la construction de la surface, il faut prendre sur la droite men\'ee par deux points quelconques des deux courbes de poursuite respectivement les points doubles (foyers) de l'involution form\'ee par ces deux points et les deux points de rencontre de la droite avec l'absolue.

\newpage
% ========================================================================
% PAPER 900: JAMES JOSEPH SYLVESTER
% ========================================================================

\section*{900. James Joseph Sylvester.}
\addcontentsline{toc}{section}{900. James Joseph Sylvester.}

\noindent\textit{[(Scientific Worthies in Nature, xxv.); Nature, vol.\ xxxix.\ (1889), pp.\ 217--219.]}

\vspace{0.5em}
\noindent\textbf{[p.\ 43]}

James Joseph Sylvester, born in London on September 3, 1814, is the sixth and youngest son of the late Abraham Joseph Sylvester, formerly of Liverpool*. He was educated at two private schools in London, and at the Royal Institution, Liverpool, whence he proceeded in due course of time to St John's College, Cambridge.

In these early days he manifested considerable aptitude for mathematics, and so it was not matter for surprise that he came out in the Tripos Examination of 1837 as Second Wrangler; being incapacitated, by the fact of his Jewish origin, from taking his degree, he was not able to compete for either of the Smith's Prizes. In more enlightened times (1872), he had the degrees of B.A.\ and M.A., by accumulation, conferred upon him, and received therewith the honour of a Latin speech from the Public Orator. He himself says: ``I am perhaps the only man in England who am a full (voting) Master of Arts for the three Universities of Dublin, Cambridge, and Oxford, having received that degree from these Universities in the order above given: from Dublin, by \textit{ad eundem}; from Cambridge, \textit{ob merita}; from Oxford, by decree.'' He is now D.C.L.\ of Oxford, LL.D.\ of Dublin and Edinburgh, and Hon.\ Fellow of St John's College, Cambridge. It is still open for him to receive yet higher recognition from his own alma mater**.

Prof.\ Sylvester became a student of the Inner Temple, July 29, 1846, and was called to the Bar on November 22, 1850\dag. He has been Professor of Natural Philosophy at University College, London; of Mathematics at the University of Virginia, U.S.A.\ddag; then ten years later Professor at the Royal Military Academy, Woolwich; and again, after a five years' interval, Professor of Mathematics at the Johns Hopkins University, Baltimore, U.S.A., from its foundation in 1877. Finally, in December 1883, he was elected Savilian Professor of Geometry at Oxford, in succession to Prof.\ Henry Smith*. His first printed paper was on Fresnel's optical theory (in the Phil.\ Mag., 1837).

\vspace{0.3em}
\noindent * Foster's \textit{Hand-book of Men at the Bar}.\\
\noindent **[The honorary degree Sc.D.\ was conferred upon him by Cambridge in 1890.]\\
\noindent \dag Foster, l.c.\\
\noindent \ddag The late Prof.\ Key, of University College and School, was the first occupant of the Chair, founded by Mr Jefferson, once President of the United States, in 1824.

\vspace{0.5em}
\noindent\textbf{[p.\ 44]}

We can here only briefly allude to a communication which was accompanied by many important results: we refer to the Friday evening address (January 23, 1874) to the Royal Institution, ``On Recent Discoveries in Mechanical Conversion of Motion.'' He says:---``It would be difficult to quote any other discovery which opens out such vast and varied horizons as this of Peaucellier's,---in one direction, descending to the wants of the workshop, the simplification of the steam-engine, the revolutionising of the mill-wright's trade, the amelioration of garden-pumps, and other domestic conveniences (the sun of science glorifies all it shines upon); and in the other, soaring to the sublimest heights of the most advanced doctrines of modern analysis, lending aid to, and throwing light from a totally unexpected quarter on the researches of such men as Abel, Riemann, Clebsch, Grassmann, and Cayley. Its head towers above the clouds, while its feet plunge into the bowels of the earth.''

The only works that Prof.\ Sylvester has published, we believe, are: (1) ``A Probationary Lecture on Geometry, delivered before the Gresham Committee and the Members of the Common Council of the City of London, December 4, 1854,'' a slight thing which had to be written and delivered at a few hours' notice; (2) ``Laws of Verse,'' 1870; (3) several short poems, sonnets, and translations, which have appeared in our columns and elsewhere.

Our notice would be incomplete without some record of the honours that have been conferred upon Dr Sylvester. He was elected a Fellow of the Royal Society on April 25, 1839; has received a Royal Medal (1860) and the Copley Medal (1880), this latter rarely awarded, we believe, to a pure mathematician. On this last occasion, Mr Spottiswoode accompanied the presentation with the words, ``His extensive and profound researches in pure mathematics, especially his contributions to the theory of invariants and covariants, to

\vspace{0.5em}
\noindent\textbf{[p.\ 45]}

which he is the chief contributor, and his discoveries in the higher theory of partitions, have given him a high and permanent position among the great mathematicians of the age.''

He is Foreign Associate of the Institute of France and of the Royal Society of Gottingen; Corresponding Member of the Institute of Bologna; \&c., \&c. He is one of the original members of the Order of Merit; and was the first recipient of the Gold Medal of the Royal Society of Sciences in Christiania. In connexion with the Royal Society, his work as chairman of the Council (to which he was elected in 1880 in succession to Mr Spottiswoode) may be referred to; he was also Chairman of the University Commission, and has been from the first one or another of its titles), and was the first editor of, and is a considerable contributor to, the American Journal of Mathematics; and he was at one time Examiner in Mathematics and Natural Philosophy in the University of London. He was not an original member of the London Mathematical Society (founded January 16, 1865), but was elected a member on June 19, 1865, Vice-President on January 15, 1866, and succeeded Prof.\ De Morgan as the second President on November 8, 1866. The Society showed its recognition of his great services to them and to mathematical science generally by awarding him its De Morgan Gold Medal in November 1887.

Wherever Dr Sylvester goes, there is sure to be mathematical activity; and the latest proof of this is the formation, during the last term at Oxford, of a Mathematical Society, which promises, we hear without surprise, to do much for the advancement of mathematical science there.

The writings of Sylvester date from the year 1837; the number of them in the Royal Society Index up to the year 1863 is 112, in the next ten years 38, and in the volume for the next ten years 81, making 231 for the years 1837 to 1883: the number of more recent papers is also considerable. They relate chiefly to finite analysis, and cover by their subjects a large part of it: algebra, determinants, elimination, the theory of equations, partitions, tactic, the theory of forms, matrices, reciprocants, the Hamiltonian numbers, \&c.; analytical and pure geometry occupy a less prominent position; and mechanics, optics, and astronomy are not absent. A leading feature is the power which is shown of originating a theory or of developing it from a small beginning; there is a breadth of treatment and determination to make the work

\vspace{0.5em}
\noindent\textbf{[p.\ 47]}

accessible, and an absence of that pride which too often refuses to expound the meaning and scope of a great discovery, and would as soon bequeath to posterity a hieroglyph as an algebraical formula.

There is, in 1844, in the Philosophical Magazine, a valuable paper, ``Elementary Researches in the Analysis of Combinatorial Aggregation,'' and the titles of two other papers, 1865 and 1866, may be mentioned: ``Astronomical Prolusions; commencing with the instantaneous proof of Lambert's and Euler's theorems, and modulating through the construction of the orbit of a heavenly body from two heliocentric distances, the subtended chord, and the periodic time, and the focal theory of Cartesian ovals, into a discussion of motion in a circle and its relation to planetary motion''; and the sequel thereto, ``Note on the periodic changes of orbit under certain circumstances of a particle acted upon by a central force, and on vectorial coordinates, \&c., together with a new theory of the analogues of the Cartesian ovals in space.''

Many of the later papers are published in the American Mathematical Journal, founded, in 1878, under the auspices of the Johns Hopkins University, and for the first six volumes of which Sylvester was editor-in-chief. We have, in vol.\ I., a somewhat speculative paper entitled ``An application of the new atomic theory to the graphical representation of the invariants and covariants of binary quantics,'' followed by appendices and notes relating to various special points of the theory; and in the same and subsequent volumes various memoirs on binary and ternary quantics, including papers (by himself, with the aid of Franklin) containing tables of the numerical generating functions for binary quantics of the first ten orders, and for simultaneous binary quantics of the first four orders, \&c. The memoir (vols.\ II.\ and III.) on ``Ternary cubic-form equations'' is connected with some early papers relating to the theory of numbers. We have in it the theory of residuation on a cubic curve, and the beautiful chain-rule of rational derivation; viz. from an arbitrary point 1 on the curve it is possible to derive the singly infinite series of points $(1, 2, 4, 5, \ldots, 3\mu \pm 1)$ such that the chord through any two points, $m$, $n$, again meets the curve in a point $m + n$, $m - n$ (whichever number is not divisible by 3) of the series; moreover, the coordinates of any point $m$ are rational and integral functions of the degree $m^2$ of those of the point 1.

There is in vol.\ v.\ the memoir, ``A Constructive Theory of Partitions arranged in three acts, an Interact in two parts, and an Exodion,'' and in vol.\ vi.\ we have ``Lectures on the Principles of Universal Algebra,'' (referring to a course of lectures on multinomial quantity, in the year 1881). The memoir is incomplete, but the general theories of nullity and vacuity, and of the corpus formed by two independent matrices of the same order, are sketched out; and there are, in the Comptes rendus of the French Academy, later papers containing developments of various points of the theory,---the conception of ``nivellators'' may be referred to.

The last-mentioned paper in the American Mathematical Journal was published subsequently to Sylvester's return to England on his appointment as Savilian Professor of Mathematics at Oxford. In December 1886, he gave there a public lecture containing an outline of his new theory of reciprocants (reported in Nature, January 7, 1887), and the lectures since delivered are published under the title, ``Lectures on the Theory of Reciprocants'' (reported by J.\ Hammond), same Journal, vols.\ VIII.\ to X.; thirty-three lectures actually delivered, entire or in abstract, in the course of

\vspace{0.5em}
\noindent\textbf{[p.\ 48]}

three terms, to a class in the University, with a concluding so-called lecture 34, which is due to Hammond. The subject, as is well known, is that of the functions of a dependent variable, $y$, and its differential coefficients, $y'$, $y''$, \ldots, in regard to $x$ (or, rather, the functions of $y'$, $y''$, \ldots), which remain unaltered by the interchange of the variables $x$ and $y$: this is a less stringent condition than that imposed by Halphen (``Th\`ese,'' 1878) on his differential invariants, and the theory is accordingly a more extensive one. A passage may be quoted:---``One is surprised to reflect on the change which is come over Algebra in the last quarter of a century. It is now possible to enlarge to an almost unlimited extent on any branch of it. These thirty lectures, embracing only a fragment of the theory of reciprocants, might be compared to an unfinished epic in thirty cantos. Does it not seem as if Algebra had attained to the dignity of a fine art, in which the workman has a free hand to develop his conceptions, as in a musical theme or a subject for painting? Formerly, it consisted in detached theorems, but nowadays it has reached a point in which every properly-developed algebraical composition, like a skilful landscape, is expected to suggest the notion of an infinite distance lying beyond the limits of the canvas.'' And, indeed, the theory has already spread itself out far and wide, not only in these lectures by its founder, but in various papers by auditors of them, and others,---Elliott, Hammond, Leudesdorf, Rogers, Macmahon, Berry, Forsyth.

Sylvester's latest important investigations relate to the Hamiltonian numbers: there is a memoir, Crelle, t.\ c.\ (1887), and, by Sylvester and Hammond jointly, two memoirs in the Philosophical Transactions. The subject is that of the series of numbers $2, 3, 5, 11, 47, 923$, calculated thus far by Sir W.\ R.\ Hamilton in his well-known Report to the British Association, on Jerrard's method. A formula for the independent calculation of any term of the series was obtained by Sylvester, but the remarkable law by means of a generating function was discovered by Hammond, viz. $E_0$, $E_1$, $E_2$, \ldots, being the series $3, 4, 6, \ldots$ of the foregoing numbers, each increased by unity; then these are calculated by the formula
$$
\frac{1 + t}{1 - t} = \prod_{n=0}^{\infty} (1 - t^{E_n}),
$$
equating the powers of $t$ on the two sides respectively: observe the paradox, $t = \frac{1}{2}$, then the formula gives $0 = \text{sum of a series of positive powers of } \frac{1}{2}$.

Enough has been said to call to mind some of Sylvester's achievements in mathematical science. Nothing further has been attempted in the foregoing very imperfect sketch.

\newpage
% ========================================================================
% PAPER 901: NOTE ON THE SUMS OF TWO SERIES
% ========================================================================

\section*{901. Note on the Sums of Two Series.}
\addcontentsline{toc}{section}{901. Note on the Sums of Two Series.}

\noindent\textit{[From the Messenger of Mathematics, vol.\ xix.\ (1890), pp.\ 29--31.]}

\vspace{0.5em}
\noindent\textbf{[p.\ 49]}

I consider the two series
$$
S_1 = \frac{1}{1 + e^a} + \frac{1}{3(1 + e^{3a})} + \frac{1}{5(1 + e^{5a})} + \cdots
$$
and
$$
S_2 = \frac{1}{1 - e^a} + \frac{1}{3(2 + 3e^a)} + \frac{1}{5(2 + 5e^a)} + \cdots,
$$
where $a$ is real, positive, and indefinitely small; these would at first sight appear to be equal to each other, but this is not in fact the case.

Taking first the series $S_1$; putting therein $\pi a = 2x$, this is
$$
S_1 = \frac{1}{1 + x} + \frac{1}{3(1 + 3x)} + \frac{1}{5(1 + 5x)} + \cdots
$$
Now we have (Legendre, \textit{Th\'eorie des Fonctions Elliptiques}, t.\ II.\ p.\ 438),
$$
\psi(y) = -C + \log y + \frac{1}{2y} - \frac{B_1}{2y^2} + \frac{B_2}{4y^4} - \cdots,
$$
where $C$ is Euler's constant, $= \cdot 577\ldots$; and if $y$ be real, positive, and very large, then
$$
\psi(y) = -C + \log y + \frac{1}{2y},
$$
whence, differentiating the logarithm and neglecting the terms which contain negative powers of $y$, then the value is $= C + \log y$; hence, writing $y = \frac{1}{x}$, we obtain
$$
\frac{1}{1 + x} + \frac{1}{2(1 + 2x)} + \frac{1}{3(1 + 3x)} + \cdots = C + \log\frac{1}{x}.
$$

\vspace{0.5em}
\noindent\textbf{[p.\ 50]}

Writing herein $2x$ for $x$, and dividing by 2, we have
$$
\frac{1}{2(1 + 2x)} + \frac{1}{4(1 + 4x)} + \frac{1}{6(1 + 6x)} + \cdots = \frac{1}{2}\left(C + \log\frac{1}{2x}\right).
$$
Hence, subtracting,
$$
S_1 = \frac{1}{2}\left(C + \log\frac{1}{x}\right) - \frac{1}{2}\left(C + \log\frac{1}{2x}\right) = \frac{1}{2}\log 2.
$$
Hence, writing for $x$ its value, $= \frac{1}{2}\pi a$, we have
$$
S_1 = \frac{1}{2}\log 2.
$$

For the series $S_2$, we have (\textit{Fundamenta Nova}, p.\ 103*) the formula
$$
\frac{2K}{\pi} = \frac{1}{1 + q^{-1}} + \frac{1}{3(1 + q^{-3})} + \frac{1}{5(1 + q^{-5})} + \cdots,
$$
where
$$
q = e^{-\pi\frac{K'}{K}}.
$$
Or, putting herein $a = \frac{K'}{K}$, then $q = e^{-\pi a} = e^{-\pi a}$, and thence
$$
S_2 = \frac{1}{1 + e^{-\pi a}} + \frac{1}{2(1 + e^{-2\pi a})} + \frac{1}{3(1 + e^{-3\pi a})} + \cdots = \frac{2K}{\pi}.
$$
We have $q = e^{-\pi a}$, which is real, positive, and less than but indefinitely near to 1; hence also $k$ is real, positive, and less than but indefinitely near to 1, say the value is $= 1 - \beta$; thence $k' = \sqrt{2\beta}$, and $K = \log\frac{4}{\sqrt{\beta}}$, $= \log\frac{4}{\sqrt{k'}}$; also $K' = \frac{1}{2}\pi$, and therefore
$$
a = \frac{K'}{K} = \frac{\frac{1}{2}\pi}{\log\frac{4}{\sqrt{k'}}}, \quad \text{whence } \log\frac{4}{\sqrt{k'}} = \frac{\pi}{2a},
$$

\newpage
% ========================================================================
% PAPER 902: ON THE FOCALS OF A QUADRIC SURFACE
% ========================================================================

\section*{902. On the Focals of a Quadric Surface.}
\addcontentsline{toc}{section}{902. On the Focals of a Quadric Surface.}

\noindent\textit{[From the Messenger of Mathematics, vol.\ xix.\ (1890), pp.\ 113--117.]}

\vspace{0.5em}
\noindent\textbf{[p.\ 51]}

In plane geometry, the focus of a curve is the node of the circumscribed line-system of the curve and the circular points at infinity; and so, in solid geometry, the focal of a surface or curve is the nodal line of the circumscribed developable of the surface or curve and the circle at infinity. And as in plane geometry the circular points at infinity may be regarded as an indefinitely thin conic, so in solid geometry the circle at infinity may be regarded as an indefinitely thin quadric surface.

In plane geometry, let it be proposed to find the circumscribed line-system (common tangents) of the two quadrics
$$
\frac{x^2}{a} + \frac{y^2}{b} + \frac{z^2}{c} = 0, \quad \frac{x^2}{a'} + \frac{y^2}{b'} + \frac{z^2}{c'} = 0;
$$
if a common tangent be
$$
\xi x + \eta y + \zeta z = 0,
$$
then we have
$$
a\xi^2 + b\eta^2 + c\zeta^2 = 0, \quad a'\xi^2 + b'\eta^2 + c'\zeta^2 = 0;
$$
here writing for shortness
$$
f, g, h = bc' - b'c, \quad ca' - c'a, \quad ab' - a'b,
$$
we have
$$
\xi^2 : \eta^2 : \zeta^2 = f : g : h;
$$
and thence the tangent is
$$
\sqrt{fx} + \sqrt{gy} + \sqrt{hz} = 0,
$$
viz.\ we have thus four tangents, and the rationalised form is of course
$$
f^2 x^4 + g^2 y^4 + h^2 z^4 - 2ghy^2 z^2 - 2hfz^2 x^2 - 2fgx^2 y^2 = 0.
$$

\vspace{0.5em}
\noindent\textbf{[p.\ 52]}

In connexion with the corresponding question in solid geometry, I obtain this equation in a different manner. We investigate the envelope of the line $\xi x + \eta y + \zeta z = 0$, considering $\xi$, $\eta$, $\zeta$ as parameters connected by the foregoing two equations; by the ordinary process of indeterminate multipliers, we have
$$
\frac{x^2}{\lambda a + \mu a'} + \frac{y^2}{\lambda b + \mu b'} + \frac{z^2}{\lambda c + \mu c'} = 0,
$$
and thence, eliminating $\xi$, $\eta$, $\zeta$, we obtain
$$
\frac{x^2}{\lambda a + \mu a'} + \frac{y^2}{\lambda b + \mu b'} + \frac{z^2}{\lambda c + \mu c'} = 0,
$$
equations equivalent to two equations, from which $\lambda$ and $\mu$ are to be eliminated. The second and third equations are the derived functions of the first equation in regard to $\lambda$, $\mu$ respectively; and hence, expressing the first equation in an integral form, the result is
$$
\Disc\, [x^2(\lambda b + \mu b')(\lambda c + \mu c') + \text{\&c.}] = 0;
$$
viz.\ this is
$$
\{(bc' + b'c)x^2 + (ca' + c'a)y^2 + (ab' + a'b)z^2\}^2 - 4(bcx^2 + cay^2 + abz^2)(b'c'x^2 + c'a'y^2 + a'b'z^2) = 0,
$$
or developing and reducing, we have the foregoing result, the coefficients entering through the combinations
$$
f, g, h = bc' - b'c, \quad ca' - c'a, \quad ab' - a'b.
$$
Writing $c = -1$, also $a = b = 1$, $c' = 0$: the equation of the first conic is
$$
\frac{x^2}{a} + \frac{y^2}{b} - z^2 = 0,
$$
and that of the second may be replaced by the two equations $x^2 + y^2 = 0$, $z = 0$, viz.\ these give the circular points at infinity: we have $f, g, h = 1, -1, a - b$, and the equation of the line-system is
$$
x^4 + 2x^2 y^2 + y^4 - 2h(x^2 - y^2)z^2 + h^2 z^4 = 0.
$$
If finally, $z = 1$, then for the tangents from the circular points at infinity to the quadric
$$
\frac{x^2}{a} + \frac{y^2}{b} - 1 = 0,
$$
the equation is
$$
x^4 + 2x^2 y^2 + y^4 - 2h(x^2 - y^2) + h^2 = 0,
$$
where, as before, $h = a - b$; the four tangents intersect in pairs in the two circular points at infinity, and in four other points which are the foci of the quadric.

\vspace{0.5em}
\noindent\textbf{[p.\ 53]}

Passing now to the problem in solid geometry, we consider the two quadric surfaces
$$
\frac{x^2}{a} + \frac{y^2}{b} + \frac{z^2}{c} - \frac{w^2}{d} = 0, \quad \frac{x^2}{a'} + \frac{y^2}{b'} + \frac{z^2}{c'} - \frac{w^2}{d'} = 0;
$$
if a common tangent plane hereof be
$$
\xi x + \eta y + \zeta z + \omega w = 0,
$$
then we have
$$
a\xi^2 + b\eta^2 + c\zeta^2 + d\omega^2 = 0, \quad a'\xi^2 + b'\eta^2 + c'\zeta^2 + d'\omega^2 = 0;
$$
and the circumscribed developable is the envelope of the plane, considering in the equation thereof $\xi$, $\eta$, $\zeta$, $\omega$ as variable parameters connected by the last two equations.

By what precedes, it is at once seen that the resulting equation is
$$
\Disc\, [x^2(b\lambda + b'\mu)(c\lambda + c'\mu)(d\lambda + d'\mu) + \text{\&c.}] = 0,
$$
viz.\ this is a quadric equation in $(x^2, y^2, z^2, w^2)$, the coefficients $a, b, c, d, a', b', c', d'$, entering therein through the combinations
$$
\alpha, \beta, \gamma = bc' - b'c, \quad ca' - c'a, \quad ab' - a'b;
$$
$$
f, g, h = ad' - a'd, \quad bd' - b'd, \quad cd' - c'd.
$$
The developed result is given in my paper ``On the developable surfaces which arise from two surfaces of the second order,'' Camb.\ and Dub.\ Math.\ Jour., t.\ V.\ (1850), pp.\ 45--53, [84]. We require here, for the particular case, $d = -1$, $a' = b' = c' = 1$, $d' = 0$, viz.\ one of the surfaces is taken to be
$$
\frac{x^2}{a} + \frac{y^2}{b} + \frac{z^2}{c} = 1,
$$
and the other is the circle at infinity $x^2 + y^2 + z^2 = 0$, $w = 0$; we thus have $\alpha, \beta, \gamma = 1, 1, 1$; $f, g, h = b - c, c - a, a - b$, or now using the italic letters $f, g, h$ to signify these values, we have $f + g + h = 0$. The equation is
\begin{multline*}
f^2 g^2 h^2 w^8 \\
+ 2g^2(g - h)y^2 z^6 + 2h^2(h - f)z^2 x^6 + 2f^2(f - g)x^2 y^6 \\
+ 2g^2(h - f)y^6 w^2 + 2h^2(f - g)z^6 w^2 + 2f^2(g - h)x^6 w^2 \\
- 2fgh^2(h - f)y^2 z^4 w^2 - 2f^2 g h(g - h)z^2 x^4 w^2 - 2fgh^2(f - g)x^2 y^4 w^2 \\
+ \cdots = 0.
\end{multline*}

\vspace{0.5em}
\noindent\textbf{[p.\ 54]}

The equation may be written in the following four forms:
\begin{align*}
x^2\Theta_1 + (gy^2 - hz^2 + ghw^2)^2\{y^4 + 2y^2 z^2 + z^4 - 2f(y^2 - z^2)w^2 + f^2 w^4\} &= 0, \\
y^2\Theta_2 + (hz^2 - fx^2 + hfw^2)^2\{z^4 + 2z^2 x^2 + x^4 - 2g(z^2 - x^2)w^2 + g^2 w^4\} &= 0, \\
z^2\Theta_3 + (fx^2 - gy^2 + fgw^2)^2\{x^4 + 2x^2 y^2 + y^4 - 2h(x^2 - y^2)w^2 + h^2 w^4\} &= 0, \\
w^2\Theta_4 + (fx^2 + gy^2 + hz^2)^2(x^2 + y^2 + z^2)^2 &= 0;
\end{align*}
the last of these shows that the circle at infinity $x^2 + y^2 + z^2 = 0$, $w = 0$, is a nodal line on the surface; this, however, is not regarded as a focal. The other three show that we have also, as nodal lines, three conics, which are the focals of the given surface; viz.\ now writing $w = 1$, we have, for the quadric surface
$$
\frac{x^2}{a} + \frac{y^2}{b} + \frac{z^2}{c} = 1,
$$
the three focal conics
\begin{align*}
x = 0, \quad \frac{y^2}{a - b} + \frac{z^2}{a - c} - 1 &= 0, \\
y = 0, \quad \frac{x^2}{a - b} + \frac{z^2}{b - c} - 1 &= 0, \\
z = 0, \quad \frac{x^2}{a - c} + \frac{y^2}{b - c} - 1 &= 0;
\end{align*}
viz.\ these are an imaginary conic, the focal hyperbola, and the focal ellipse, in the coordinate planes $x = 0$, $y = 0$, $z = 0$ respectively.

\newpage
% ========================================================================
% PAPER 903: ON LATIN SQUARES
% ========================================================================

\section*{903. On Latin Squares.}
\addcontentsline{toc}{section}{903. On Latin Squares.}

\noindent\textit{[From the Messenger of Mathematics, vol.\ xix.\ (1890), pp.\ 135--137.]}

\vspace{0.5em}
\noindent\textbf{[p.\ 55]}

If in each line of a square of $n^2$ compartments the same $n$ letters $a, b, c, \ldots$ are arranged so that no letter occurs twice in the same column, we have what was termed by Euler ``a Latin square.'' Supposing that in one of the lines the letters are arranged in the natural order $abcde\ldots$, then in the remaining lines there must be arrangements beginning with $b, c, d, e, \&c.$, respectively, and we may consider the case in which the bottom line has the arrangement $abcde\ldots$, and in the other lines, reckoning from the bottom one in order, the arrangements begin with $b, c, d, e, \&c.$, respectively: if the number of such squares be $= N$, then, obviously, the whole number of squares which can be formed with the same $n$ arrangements is $= N[n]^n$.

Starting with the bottom line as above, then it is a well-known problem to determine the number of arrangements for the second line; viz.\ this number is
$$
[n]^n - \frac{n}{1}[n-1]^{n-1} + \frac{n(n-1)}{1 \cdot 2}[n-2]^{n-2} - \cdots \pm 1,
$$
and if we assume, as above, that the second line begins with $b$, then the whole number of arrangements is this number divided by $(n - 1)$, the quotient being of course integral. For instance, when $n = 5$, the number is $= 120 - 120 + 60 - 20 + 5 - 1 = 44$, which is divisible by 4, and the number of arrangements for the second line is thus $= 11$.

But the number of arrangements for the third line will be different according to the arrangement selected for the second line, and it is not easy to see how in general the whole number of arrangements for the third line is to be calculated, and the difficulty of course increases for the next following lines; it may be remarked that, when all the lines are filled up except the top-line, the top-line is completely determined.

\vspace{0.5em}
\noindent\textbf{[p.\ 56]}

Imagine the square completed: we may write down the substitutions by which we pass from the bottom line to itself (this is of course the substitution 1) and to each of the other lines respectively; we have thus a set of $n$ substitutions, which may form a group; and when this is so, we may conversely from the group construct the Latin square. But it is not every Latin square which is thus connected with a group of $n$ substitutions.

In the cases $n = 2, 3, 4$ there is no difficulty; the squares are
$$
\begin{array}{cc}
b & a \\
a & b
\end{array},
\quad
\begin{array}{ccc}
c & a & b \\
b & c & a \\
a & b & c
\end{array},
\quad
\begin{array}{cccc}
d & c & b & a \\
c & d & a & b \\
b & a & d & c \\
a & b & c & d
\end{array},
\quad
\begin{array}{cccc}
d & c & b & a \\
c & d & a & b \\
b & d & a & c \\
a & b & c & d
\end{array},
\quad
\begin{array}{cccc}
d & a & b & c \\
c & a & d & b \\
b & c & d & a \\
a & b & c & d
\end{array},
\quad
\begin{array}{cccc}
d & a & b & c \\
c & d & a & b \\
b & c & d & a \\
a & b & c & d
\end{array};
$$
viz.\ when $n = 2$ the number is 1, when $n = 3$ it is 1, when $n = 4$ it is 4; in this last case, the arrangement $bcda$ for the second line gives two squares, but each of the other arrangements only one square.

In each of the squares of 4, we have a group, viz.\ for the four squares respectively, these are
\begin{align*}
1^\circ &\quad 1, \quad (ab)(cd), \quad (acbd), \quad (adbc), \\
2^\circ &\quad 1, \quad (ab)(cd), \quad (ac)(bd), \quad (ad)(bc), \\
3^\circ &\quad 1, \quad (abdc), \quad (acdb), \quad (ad)(bc), \\
4^\circ &\quad 1, \quad (abcd), \quad (ad)(bc), \quad (ac)(bd).
\end{align*}

\vspace{0.5em}
\noindent\textbf{[p.\ 57]}

The case $n = 5$ is much more difficult; I have not as yet completely solved it, but I have obtained the following results. Denoting the arrangements for the second line by $2, 3, 4, 5$ (each of these denotes an arrangement beginning with the corresponding letter), I find that 2 gives 11 squares, 3 gives 11, 4 gives 11, and 5 gives 2; the whole number of squares is thus $= 35$, agreeing with Prof.\ Cayley's former determination. The groups of 5 are very remarkable; I have not investigated them completely, but there are at least 5 different types of group which occur; viz.\ one group of the cyclical type, and four or more groups of other types.

The cyclical group belongs to the squares corresponding to the arrangement 5 of the second line; it is generated by the substitution $(abcde)$, and the corresponding squares are
$$
\begin{array}{ccccc}
a & b & c & d & e \\
b & c & d & e & a \\
c & d & e & a & b \\
d & e & a & b & c \\
e & a & b & c & d
\end{array}
\quad \text{and} \quad
\begin{array}{ccccc}
d & e & a & b & c \\
c & d & e & a & b \\
b & c & d & e & a \\
a & b & c & d & e
\end{array}.
$$
These are the only squares belonging to the cyclical group; but each of the same arrangement 5 gives other squares belonging to non-cyclical groups.

\newpage
% ========================================================================
% PAPER 904: NOTE ON RECIPROCAL LINES
% ========================================================================

\section*{904. Note on Reciprocal Lines.}
\addcontentsline{toc}{section}{904. Note on Reciprocal Lines.}

\noindent\textit{[From the Messenger of Mathematics, vol.\ xix.\ (1890), pp.\ 174, 175.]}

\vspace{0.5em}
\noindent\textbf{[p.\ 58]}

If two lines are reciprocal in regard to a quadric surface, then any point on the one line and any point on the other line are harmonics in regard to the surface, viz.\ the two points and the intersections of their line of junction with the surface form a harmonic range. This is obvious: the polar plane of the first point passes through the second line, and thus the second point is a point in the polar plane of the first point, that is, the two points are harmonics.

But it is worth while to look at the theorem in a different point of view. If the first line meets the surface in the points $A$, $C$ and the second line meets the surface in the points $B$, $D$ (in order to the four points being real, the surface must of course be a skew hyperboloid), then $AB$, $BC$, $CD$, $DA$ are lines on the surface, say $AB$, $CD$ are directrices, and $BC$, $AD$ are generatrices, the two reciprocal lines being the diagonals $AC$ and $BD$ of the skew quadrilateral. Taking on $AC$ a point $P$ and on $BD$ a point $Q$, the theorem is that, if $PQ$ meets the surface in the points $X$, $Y$, then the four points $P$, $Q$, $X$, $Y$ are harmonics. Let the generatrix through $X$ meet $AB$, $CD$ in the points $S$, $S'$ respectively, and the generatrix through $Y$ meet the same lines in the points $T$, $T'$ respectively.

\vspace{0.5em}
\noindent\textbf{[p.\ 59]}

Then the four lines $CB$, $DA$, $T'T$, $S'S$, are all met by an infinite series of lines, and in particular they are met by the lines $CD$, $BA$ in the points $C$, $D$, $T'$, $S'$ and $B$, $A$, $T$, $S$ respectively. Consequently, writing AH for anharmonic ratio, we have
$$
\text{AH}(C, D, T', S') = \text{AH}(B, A, T, S).
$$
Consider now the four lines $CA$, $DB$, $T'T$, $S'S$: these are each of them met by the lines $CD$ and $BA$, and also by the line $PQ$: therefore, by an infinity of lines (that is, they are directrices of a hyperboloid). They are met by the last-mentioned three lines in the points $C$, $D$, $T'$, $S'$; $A$, $B$, $T$, $S$; and $P$, $Q$, $Y$, $X$ respectively; hence
$$
\text{AH}(C, D, T', S') = \text{AH}(A, B, T, S) = \text{AH}(P, Q, Y, X).
$$
Comparing with the foregoing equation, it appears that
$$
\text{AH}(B, A, T, S) = \text{AH}(A, B, T, S),
$$
viz.\ the anharmonic ratio is not altered by the interchange of the two points $A$, $B$: this implies that the points $A$, $B$, $T$, $S$ are harmonics, viz.\ the pairs $A$, $B$ and $S$, $T$ are harmonics; and, this being so, the equation
$$
\text{AH}(A, B, T, S) = \text{AH}(P, Q, Y, X)
$$
shows that $P$, $Q$, $Y$, $X$ are harmonics, viz.\ that the pairs $P$, $Q$ and $X$, $Y$ are harmonics; or say $P$, $Q$ are harmonics in regard to the quadric surface, which is the theorem in question.

\newpage
% ========================================================================
% PAPER 905: ON THE EQUATION $x^{17} - 1 = 0$
% ========================================================================

\section*{905. On the Equation $x^{17} - 1 = 0$.}
\addcontentsline{toc}{section}{905. On the Equation $x^{17} - 1 = 0$.}

\noindent\textit{[From the Messenger of Mathematics, vol.\ xix.\ (1890), pp.\ 184--188.]}

\vspace{0.5em}
\noindent\textbf{[p.\ 60]}

Writing $\rho = \cos\frac{2\pi}{17} + i\sin\frac{2\pi}{17}$, I carry the solution up to the determination of the periods each of two roots, $\rho + \rho^{16}$, $= 2\cos\frac{2\pi}{17}$, \&c. The expressions contain the radicals
$$
a = \sqrt{17}, \quad b = \sqrt{\{2(17 - a)\}}, \quad c = \sqrt{\{4(17 + 3a) - 2(3 + a)b\}},
$$
where $a$, $b$, $c$ are taken to be positive ($a = 4\cdot 12$, $b = 5\cdot 07$, $c = 6\cdot 72$). Taking for a moment $r$ to be any imaginary seventeenth root, $r = \rho^\theta$, then the algebraical expression for the period $P_1$ of eight roots is $P_1 = \frac{1}{2}(-1 + a)$, but I assume the value to be $-P_1 = \frac{1}{2}(-1 + a)$, and thus determine $\theta$ to denote some one of the values $1, 2, 4, 8, 9, 13, 15, 16$; similarly, I assume the value of $Q_1$ to be $= \frac{1}{2}(-1 + a + b)$; and thus further determine $\theta$ to denote some one of the values $1, 4, 13, 16$: and, again, I assume the value of $R_1$ to be $= \frac{1}{2}(-1 + a + b + c)$, and thus further determine $\theta$ to denote one of the values 1 and 16. As regards the values of the periods $R$, it is obviously indifferent which value is taken, and I assume therefore $\theta = 1$. This comes to saying that the signs of the radicals are determined in suchwise that $r$ shall denote the root $\cos\frac{2\pi}{17} + i\sin\frac{2\pi}{17}$; and it is to be understood that $r$ has this value.

I write now
\begin{align*}
P_1 &= r + r^2 + r^4 + r^8 + r^9 + r^{13} + r^{15} + r^{16}, \\
P_2 &= r^3 + r^6 + r^5 + r^{11} + r^{14} + r^7 + r^{12} + r^6, \\
Q_1 &= r + r^4 + r^{13} + r^{16}, \\
Q_2 &= r^2 + r^8 + r^9 + r^{15}, \\
Q_3 &= r^3 + r^5 + r^{12} + r^{14}, \\
Q_4 &= r^6 + r^7 + r^{10} + r^{11},
\end{align*}
and moreover
\begin{align*}
R_1 &= r + r^{16}, \quad R_2 = r^4 + r^{13}, \quad R_3 = r^2 + r^{15}, \quad R_4 = r^8 + r^9, \\
R_5 &= r^3 + r^{14}, \quad R_6 = r^5 + r^{12}, \quad R_7 = r^{10} + r^7, \quad R_8 = r^{11} + r^6, \\
a &= P_1 - P_2, \\
b^2 &= 2(17 - a), \quad b_1^2 = 2(17 + a), \\
c^2 &= 4(17 + 3a) - 2(3 + a)b.
\end{align*}

\vspace{0.5em}
\noindent\textbf{[p.\ 61]}

It will appear by what follows that $a$ is determined by the quadric equation $a^2 = 17$, but I have assumed that $a$ denotes the positive root $a = \sqrt{17}$; similarly, $b$ is determined by the quadric equation $b^2 = 2(17 - a)$, but it is assumed that $b$ denotes the positive root, $b = \sqrt{\{2(17 - a)\}}$; and $c$ is determined by the quadric equation $c^2 = 4(17 + 3a) - 2(3 + a)b$, but it is assumed that $c$ denotes the positive root.

If in the equations I had written $b_1, c_1, c_2, c_3$, instead of $+b_1, -c_1, +c_2, -c_3$, then $b_1$ comes out rationally in terms of $a, b$; and $c_1, c_2, c_3$ come out rationally in terms of $a, b, c$; the signs were attached to them a posteriori, in suchwise that the values of $b_1, c_1, c_2, c_3$ might be each of them positive; for their independent determination, we have, in fact, for $b_1$, an expression such as that for $b_2$; and for $c_1^2, c_2^2, c_3^2$ expressions such as that for $c^2$; and taking as above for each of them the positive value of the square root, we have
\begin{align*}
b &= \sqrt{\{2(17 - a)\}} && (= 5\cdot 07), \\
b_1 &= \sqrt{\{2(17 + a)\}} && (= 6\cdot 49), \\
c &= \sqrt{\{4(17 + 3a) - 2(3 + a)b\}} && (= 6\cdot 72), \\
c_1 &= \sqrt{\{4(17 + 3a) + 2(3 + a)b_1\}} && (= 13\cdot 77), \\
c_2 &= \sqrt{\{4(17 - 3a) + 2(-3 + a)b_1\}} && (= 5\cdot 75), \\
c_3 &= \sqrt{\{4(17 - 3a) - 2(-3 + a)b_1\}} && (= 2\cdot 02).
\end{align*}

\vspace{0.5em}
\noindent\textbf{[p.\ 62]}

The relations between the periods $P$ are
$$
P_1 + P_2 = -1,
$$
that is,
$$
P_1^2 = -P_1 + 4, \quad \&c.; \quad a^2 = 17.
$$
Here for the second square, we have
$$
Q_1 + Q_2 = P_1, \quad Q_3 + Q_4 = P_2,
$$
and similarly in the other cases which follow.

For the periods $Q$, we have
\begin{align*}
Q_1^2 &= -P_1 + 4, \quad Q_1 Q_2 = P_1 + 4, \quad Q_2^2 = -P_1 + 4, \\
Q_3^2 &= -P_2 + 4, \quad Q_3 Q_4 = P_2 + 4, \quad Q_4^2 = -P_2 + 4;
\end{align*}
$= -P_1 - P_2 + 16 = 17$; so that $bb_1 = 8a$, or $b_1$ is given rationally in terms of $a, b$.

And for the periods $R$, we have
\begin{align*}
-R_1 + R_2 &= Q_1, \quad R_3 + R_4 = Q_2, \quad R_5 + R_6 = Q_3, \quad R_7 + R_8 = Q_4, \\
R_1^2 &= -R_1 + \cdots, \quad R_2^2 = -R_2 + \cdots, \quad \&c.
\end{align*}

\vspace{0.5em}
\noindent\textbf{[p.\ 63]}

And from the foregoing results, we have
\begin{align*}
P_1 &= \tfrac{1}{2}(-1 + a), \\
P_2 &= \tfrac{1}{2}(-1 - a), \\
Q_1 &= \tfrac{1}{2}(-1 + a + b), \quad = 1\cdot 87, \\
Q_2 &= \tfrac{1}{2}(-1 + a - b), \quad = 0\cdot 18, \\
Q_3 &= \tfrac{1}{2}(-1 - a + b_1), \quad = -1\cdot 96, \\
Q_4 &= \tfrac{1}{2}(-1 - a - b_1), \quad = -2\cdot 90, \\
R_1 &= \tfrac{1}{2}(-1 + a + b + c), \quad = 2\cdot 05, \\
R_2 &= \tfrac{1}{2}(-1 + a + b - c), \quad = 0\cdot 34, \\
R_3 &= \tfrac{1}{2}(-1 + a - b + c_1), \quad = 0\cdot 89, \\
R_4 &= \tfrac{1}{2}(-1 + a - b - c_1), \quad = -0\cdot 55, \\
R_5 &= \tfrac{1}{2}(-1 - a + b_1 + c_2), \quad = -1\cdot 20, \\
R_6 &= \tfrac{1}{2}(-1 - a + b_1 - c_2), \quad = -2\cdot 70, \\
R_7 &= \tfrac{1}{2}(-1 - a - b_1 + c_3), \quad = -0\cdot 49, \\
R_8 &= \tfrac{1}{2}(-1 - a - b_1 - c_3), \quad = -3\cdot 51.
\end{align*}
The approximate numerical values have been given throughout only for the purpose of showing that the signs of the square roots have been rightly determined.

\newpage
% ========================================================================
% PAPER 906: NOTE ON SCHLAEFLI'S MODULAR EQUATION FOR THE CUBIC TRANSFORMATION
% ========================================================================

\section*{906. Note on Schlaefli's Modular Equation for the Cubic Transformation; with a Correction.}
\addcontentsline{toc}{section}{906. Note on Schlaefli's Modular Equation for the Cubic Transformation.}

\noindent\textit{[From the Messenger of Mathematics, vol.\ xx.\ (1891), pp.\ 59, 60; 120.]}

\vspace{0.5em}
\noindent\textbf{[p.\ 64]}

The equation in question, Crelle, t.\ LXXII.\ (1870), p.\ 369, is
$$
S^8 + T^8 - 8ST + 8S^3 T^3 = 0,
$$
where
$$
S = \frac{\sqrt[4]{2}}{\sqrt[4]{(kk')}}, \quad T = \frac{\sqrt[4]{2}}{\sqrt[4]{(\lambda\lambda')}},
$$
which must be of course equivalent to the ordinary modular equation
$$
u^4(1 - u^8)^2 + u^4(1 - v^8)^2 - 2u^4 v^4 = 0,
$$
where
$$
u = \sqrt[4]{k}, \quad v = \sqrt[4]{\lambda};
$$
the resemblance in form between the two equations, with such different meanings of the $S$ and $T$ in the one case, and the $u$ and $v$ in the other, is very noticeable.

Schlaefli's equation is
$$
S^8 + T^8 = 8ST - 8S^3 T^3,
$$
that is,
$$
kk' + \lambda\lambda' = 2(kk')^{\frac{1}{4}} - 4(kk')^{\frac{3}{4}},
$$
or say
$$
x'^2 = 4(kk')^{\frac{1}{4}} - 18(kk')^{\frac{3}{4}} + 16(kk')^{\frac{5}{4}},
$$
where $x' = \sqrt{kk'} + \sqrt{\lambda\lambda'}$.

To deduce this from the $uv$-modular equation, we have (Jacobi's \textit{Fund.\ Nova}, p.\ 68, Ges.\ Werke, t.\ I., p.\ 124),
$$
(1 - u^8)(1 - v^8) = (1 - u^2 v^2)^4,
$$

\vspace{0.5em}
\noindent\textbf{[p.\ 65]}

or, multiplying each side by $u^4 v^4$, and extracting the fourth root, we have
$$
\sqrt{(kk'\lambda\lambda')} = uv\sqrt{(1 - u^2 v^2)} = x - x^3,
$$
if for shortness we write $x = uv$.

The equation to be proved thus is
$$
u^8(1 - u^8) + v^8(1 - v^8) = 4(x - x^3) - 18(x - x^3)^2 + 16(x - x^3)^3.
$$
But from the foregoing equation
$$
(1 - u^8)(1 - v^8) = (1 - u^2 v^2)^4,
$$
we have
$$
u^8 + v^8 = 4u^2 v^2 - 6u^4 v^4,
$$
and thence
$$
u^{16} + v^{16} = (4u^2 v^2 - 6u^4 v^4)^2 - 2u^8 v^8,
$$
and the equation to be proved thus becomes
$$
2x^4 = 4(x - x^3) - 18(x - x^3)^2 + 16(x - x^3)^3,
$$
which is in fact an identity, each side being
$$
52x^3 - 68x^5 + 16x^6.
$$

\vspace{1em}
\noindent\textbf{Correction, p.\ 120.}

I find that I misquoted Schlaefli's equation, viz.\ in effect, I took it to be
$$
S_1^8 + T_1^8 - 8S_1 T_1 + 8S_1^3 T_1^3 = 0, \quad \text{where } S_1 = \frac{2}{\sqrt[4]{(kk')}}, \quad T_1 = \frac{2}{\sqrt[4]{(\lambda\lambda')}},
$$
whereas his equation really is
$$
S^8 + T^8 - 8ST + 8S^3 T^3 = 0, \quad \text{where } S = 2\sqrt[4]{(kk')}, \quad T = 2\sqrt[4]{(\lambda\lambda')}.
$$
The change is only a change of form, for writing $S_1 = \frac{2}{S}$ and $T_1 = \frac{2}{T}$, the equation in $(S_1, T_1)$ is converted into that in $(S, T)$; but it was quite an unnecessary one, and I cannot account for having made it, as the paper in Crelle must have been before me.

\newpage
% ========================================================================
% PAPER 907: NOTE ON THE NINTH ROOTS OF UNITY
% ========================================================================

\section*{907. Note on the Ninth Roots of Unity.}
\addcontentsline{toc}{section}{907. Note on the Ninth Roots of Unity.}

\noindent\textit{[From the Messenger of Mathematics, vol.\ xx.\ (1891), p.\ 63.]}

\vspace{0.5em}
\noindent\textbf{[p.\ 66]}

Let $\theta$ be a prime ninth root of unity, so that $\theta^6 + \theta^3 + 1 = 0$; and write
$$
a = \theta + \theta^8, \quad b = \theta^2 + \theta^7, \quad c = \theta^4 + \theta^5,
$$
then
$$
a + b + c = -1,
$$
that is,
$$
ab = c + 2, \quad ca = c - 1,
$$
whence
$$
ab + ac + bc = -3, \quad abc = -1;
$$
and $a, b, c$ are thus the roots of the equation $x^3 - 3x + 1 = 0$. We have
$$
a^2 b + b^2 c + c^2 a = a^2 + b^2 + c^2 = 6,
$$
$$
ab^2 + bc^2 + ca^2 = bc + ca + ab = -3,
$$
and thence
$$
(a^2 b + b^2 c + c^2 a) - (ab^2 + bc^2 + ca^2) = (a - b)(b - c)(c - a) = 6 - (-3) = 9.
$$
The equation $x^3 - 3x + 1 = 0$ is thus such that $a^2 b + b^2 c + c^2 a$, and consequently any rational function whatever of $a, b, c$, invariable by the cyclical interchange $(abc)$ of the roots, has a rational value.

\newpage
% ========================================================================
% PAPER 908: ON TWO INVARIANTS OF A QUADRIQUADRIC FUNCTION
% ========================================================================

\section*{908. On Two Invariants of a Quadriquadric Function.}
\addcontentsline{toc}{section}{908. On Two Invariants of a Quadriquadric Function.}

\noindent\textit{[From the Messenger of Mathematics, vol.\ xx.\ (1891), pp.\ 68, 69.]}

\vspace{0.5em}
\noindent\textbf{[p.\ 67]}

The quadriquadric function
$$
z^2(ax^2 + 2hxy + gy^2) + 2zw(h'x^2 + 2bxy + fy^2) + w^2(g'x^2 + 2f'xy + cy^2),
$$
considered successively as a function of $(z, w)$ and of $(x, y)$, has the discriminants $U$, $V$, equal to
$$
(ax^2 + 2hxy + gy^2)(g'x^2 + 2f'xy + cy^2) - (h'x^2 + 2bxy + fy^2)^2,
$$
$$
(az^2 + 2h'zw + gw^2)(g'z^2 + 2f'zw + cw^2) - (hz^2 + 2bzw + f'w^2)^2,
$$
respectively. As is well known, these quartic functions have each of them the same quadrinvariant and the same cubinvariant; these are the invariants in question of the quadriquadric function.

The quadrinvariant has been calculated in a different notation, but I am not aware that the cubinvariant has been before calculated; the two values are as follows:

\vspace{0.5em}
\noindent\textbf{[p.\ 68]}

\begin{center}
\textbf{Quadrinvariant is} \qquad \textbf{Cubinvariant is}
\end{center}

\begin{center}
\begin{tabular}{r@{\hspace{2em}}r}
$a^2 c^2$ $+1$ & $a^3 c^3$ $-1$ \\
$ac(f^2 + f'^2)$ $+33$ & $acf^2 f' g h$ $-36$ \\
$acg g'$ $+33$ & $a f^3 f' g h$ $-36$ \\
$g^3 g'^3$ $-1$ & $af'^3 g' h$ $-36$ \\
$ab^4$ $-24$ & $ab^2 c g g'$ $-120$ \\
$a^2 b^2 c$ $+12$ & $a^2 c f'^2 g'$ $-36$ \\
$ab^2 c g g'$ $-120$ & $ac f' g g' h$ $-60$ \\
$a^2 c f'^2 g'$ $-36$ & $af'^2 g' h'$ $-36$ \\
$ac f' g g' h$ $-60$ & $f^2 g^2 g'^2$ $+6$ \\
$af'^2 g' h'$ $-36$ & $b^6$ $+48$ \\
$f^2 g^2 g'^2$ $+6$ & $ab^4 c$ $-48$ \\
$b^6$ $+48$ & $a^2 b^2 c f h'$ $+6$ \\
$ab^4 c$ $-48$ & $ac^2 g h^2$ $-36$ \\
$a^2 b^2 c f h'$ $+6$ & $ac f g g' h'$ $-60$ \\
$ac^2 g h^2$ $-36$ & $af' h^3$ $-36$ \\
$ac f g g' h'$ $-60$ & $f g^2 g'^2 h$ $+6$ \\
$af' h^3$ $-36$ & $ac g g'$ $+42$ \\
$f g^2 g'^2 h$ $+6$ & $b^4$ $+64$ \\
$ac g g'$ $+42$ & $a^2 c^2 f' h$ $+6$ \\
$b^4$ $+64$ & $ac^2 g' h'^2$ $-36$ \\
$a^2 c^2 f' h$ $+6$ & $ac f^2 h'^2$ $+24$ \\
$ac^2 g' h'^2$ $-36$ & $af'^2 g g'^2$ $-36$ \\
$ac f^2 h'^2$ $+24$ & $f^2 g'^2 h^2$ $+24$ \\
$af'^2 g g'^2$ $-36$ & $ac h f$ $-12$ \\
$f^2 g'^2 h^2$ $+24$ & $b^2 c^2$ $+76$ \\
$ac h f$ $-12$ & $\cdots$ \\
$b^2 c^2$ $+76$ & \\
\end{tabular}
\end{center}

[The complete tables of the quadrinvariant and cubinvariant are given in the original text.]

\newpage
% ========================================================================
% PAPER 909: ON A PARTICULAR CASE OF KUMMER'S DIFFERENTIAL EQUATION OF THE THIRD ORDER
% ========================================================================

\section*{909. On a Particular Case of Kummer's Differential Equation of the Third Order.}
\addcontentsline{toc}{section}{909. On a Particular Case of Kummer's Differential Equation of the Third Order.}

\noindent\textit{[From the Messenger of Mathematics, vol.\ xx.\ (1891), pp.\ 75--79.]}

\vspace{0.5em}
\noindent\textbf{[p.\ 69]}

The general form of equation in question is
$$
\frac{d^3 x}{dx^3} = A \frac{x'}{x^3} + B \frac{x' x''}{x^2} + C \frac{x'^3}{x^3} + A' \frac{x'}{(x - 1)^3} + B' \frac{x' x''}{(x - 1)^2} + C' \frac{x'^3}{(x - 1)^3};
$$
here $x$ is a function of $t$; and $A, B, C, A', B', C'$ are numerical constants. For various given values of $A, B, C$, and values determined thereby of $A', B', C'$, the equation admits of a solution in the form $x = \text{rational function of } t$; the theory in reference to the cases considered by Schwarz is considered in my paper ``On the Schwarzian Derivative and the Polyhedral Functions,'' Camb.\ Phil.\ Trans., t.\ XIII.\ (1883), pp.\ 5--68, [744]. But the theory is considered in a more general and exhaustive manner in Goursat's memoir, ``Recherches sur l'\'equation de Kummer,'' \textit{Acta Soc.\ Sci.\ Fennic\ae}, t.\ XV.\ (1888), pp.\ 47--127. I consider here one of the solutions given by him, viz.\ writing
\begin{align*}
P &= 4t^2 - 5t + 1, \quad X = t^2 P^3, \\
Q &= 5t - 4, \qquad\quad Y = Q^3, \\
R &= 8t^2 - 11t + 8, \quad Z = -(t - 1)^2 R,
\end{align*}
so that, identically, $X + Y + Z = 0$; then the solution is expressed by either of the equivalent equations
$$
x = \frac{X}{Z} = -\frac{t^2 P^3}{(t - 1)^2 R}, \quad \text{or} \quad 1 - x = \frac{Y}{Z} = -\frac{Q^3}{(t - 1)^2 R}.
$$

\vspace{0.5em}
\noindent\textbf{[p.\ 70]}

The values of the constants to which this solution belongs are
$$
A = -4, \quad B = 5, \quad C = -2, \quad A' = -4, \quad B' = 5, \quad C' = -2.
$$
But instead of assuming these values in the first instance, I leave the values indeterminate; and starting from the foregoing expression for $x$, I substitute this in
$$
\Omega = \frac{x'''}{x'} - 3\left(\frac{x''}{x'}\right)^2 - A \frac{x'}{x^2} - B \frac{x''}{x} - C \frac{x'^2}{x^2} - A' \frac{x'}{(x - 1)^2} - B' \frac{x''}{x - 1} - C' \frac{x'^2}{(x - 1)^2},
$$
thus obtaining $\Omega$ as a function of $t$ which, as will appear, vanishes identically when $A, B, C, A', B', C'$ have the foregoing values.

I remark that this is, in effect, doing in a somewhat different form for the particular case what Goursat does for the general case, viz.\ starting from
$$
\Omega = \frac{x'''}{x'} - 3\left(\frac{x''}{x'}\right)^2 - A \frac{x'}{x^2} - \cdots,
$$
with values of $A, B, C$ which belong to the solution considered, he shows that this is a function of $t$ having no infinities other than $(0, 1, \infty)$; that $\infty$ is not an infinity of the function or of the function multiplied into $t$, and that $0$ and $1$ are each of them a twofold infinity; that is, that the function is of the form
$$
\frac{Lt^2 + Mt + N}{t^2(t - 1)^2} - \frac{A'}{t^2} - \frac{B'}{t} - \frac{C'}{(t - 1)^2}.
$$
Proceeding to carry out the process, we have
$$
x = -\frac{t^2 P^3}{(t - 1)^2 R}, \quad x - 1 = \frac{Q^3}{(t - 1)^2 R},
$$
$$
\frac{x'}{x} = \frac{2}{t} + \frac{3P'}{P} - \frac{2}{t - 1} - \frac{R'}{R},
$$
and from either of these equations, collecting and reducing,
$$
\frac{x'}{x} = \frac{2(t - 1)P^2 \cdot tP}{t(t - 1)R} + \cdots = \frac{2(t - 1)P^2 \cdot Q^2}{t(t - 1)R},
$$
where observe that, from the values of $x$ and $x - 1$ respectively, it appears a priori that $tP^2$ and $Q^2$ must be factors in the numerator of $x'$. From this value of $x'$, we have
$$
\frac{x''}{x'} = \frac{2}{t} + \frac{2}{t - 1} + \frac{3P'}{P} + \frac{2Q'}{Q} - \frac{R'}{R},
$$
and hence, $P'$ and $Q'$ being mere constants,
$$
\frac{x'''}{x'} - 3\left(\frac{x''}{x'}\right)^2 = \cdots.
$$

\vspace{0.5em}
\noindent\textbf{[p.\ 71]}

By decomposing $\alpha\beta$, $\alpha p$, \&c., into simple fractions, this becomes
\begin{multline*}
A''\left(-4\alpha + 3\beta + \frac{4p}{t}\right) + B''\left(4\alpha - 4\beta + \frac{5q}{t - 1}\right) + C''\left(-\beta + \frac{8r}{(t - 1)^2}\right) \\
+ F'(-4\alpha + 4p) + G'(5\beta - 5q) + H'(-8r),
\end{multline*}
where $\alpha = \frac{1}{t}$, $\beta = \frac{1}{t - 1}$, $p = \frac{P'}{P}$, $q = \frac{Q'}{Q}$, $r = \frac{R'}{R}$.

This is
\begin{multline*}
+ p^2 A'' + q^2 B'' + r^2 C'' \\
+ p(-4A'' + 4F') + q(5B'' - 5G') + r(-8C'' + 8H') \\
+ \alpha(-4A'' + 4B'') + \beta(3A'' - 4B'' - C'') + \cdots.
\end{multline*}
This should be identically $= 0$; making $A'' = 0$, $B'' = 0$, $C'' = 0$, we find
$$
F' = A'', \quad G' = B'', \quad H' = C''.
$$

\vspace{0.5em}
\noindent\textbf{[p.\ 73]}

and thence
\begin{align*}
A &= -4, \quad B = 5, \quad C = -2, \\
A' &= -4, \quad B' = 5, \quad C' = -2.
\end{align*}
These values make the coefficients of $p$ and $q$ to be each $= 0$; and they make the coefficient of $R$ to be identically $= 0$, viz.\ we have
$$
0 = \tfrac{1}{2}H - 3H' - \tfrac{3}{4}F'' - 3\tfrac{1}{4}G'',
$$
and
$$
0 = \tfrac{1}{2}II - 3H' - \tfrac{3}{4}F'' - 3\tfrac{1}{4}G''.
$$
We have, moreover,
$$
L = \tfrac{1}{2} - C', \quad M = \tfrac{1}{2} - B', \quad N = 1 - A',
$$
and the coefficients of $\alpha$ and $\beta$ are $= -M + \frac{1}{2}$ and $M - \frac{1}{2}$ respectively; hence the coefficients of $\alpha^2$, $\beta^2$, $\alpha$ and $\beta$ will all vanish if only $L = 0$, $M = \frac{1}{2}$, $N = 0$, that is,
$$
A' = \frac{3}{2}, \quad B' = \frac{13}{12}, \quad C' = \frac{5}{9},
$$
and we have thus identically $\Omega = 0$, if only $A, B, C, A', B', C'$ have the above-mentioned values.

\newpage
% ========================================================================
% PAPER 910: NOTE ON THE INVOLUTANT OF TWO BINARY MATRICES
% ========================================================================

\section*{910. Note on the Involutant of Two Binary Matrices.}
\addcontentsline{toc}{section}{910. Note on the Involutant of Two Binary Matrices.}

\noindent\textit{[From the Messenger of Mathematics, vol.\ xx.\ (1891), pp.\ 136, 137.]}

\vspace{0.5em}
\noindent\textbf{[p.\ 74]}

Consider the two matrices
$$
M = \begin{pmatrix} a & b \\ c & d \end{pmatrix}, \quad M' = \begin{pmatrix} a' & b' \\ c' & d' \end{pmatrix},
$$
and their product in one or the other order
$$
MM' = \begin{pmatrix} A & B \\ C & D \end{pmatrix}, \quad M'M = \begin{pmatrix} A_1 & B_1 \\ C_1 & D_1 \end{pmatrix}.
$$
Then the Involutant is by definition $=$ either of the determinants
$$
I = \begin{vmatrix} 1 & a & a' & A \\ 0 & b & b' & B \\ 0 & c & c' & C \\ 1 & d & d' & D \end{vmatrix}, \quad
I_1 = \begin{vmatrix} 1 & a' & a & A_1 \\ 0 & b' & b & B_1 \\ 0 & c' & c & C_1 \\ 1 & d' & d & D_1 \end{vmatrix};
$$
viz.\ it is to be shown that these two values are in fact equal.

We have
$$
MM' = \begin{pmatrix} a & b \\ c & d \end{pmatrix} \begin{pmatrix} a' & b' \\ c' & d' \end{pmatrix},
$$
that is,
$$
A = aa' + bc', \quad B = ab' + bd', \quad C = ca' + dc', \quad D = cb' + dd'.
$$

\vspace{0.5em}
\noindent\textbf{[p.\ 75]}

and similarly
$$
M'M = \begin{pmatrix} a' & b' \\ c' & d' \end{pmatrix} \begin{pmatrix} a & b \\ c & d \end{pmatrix},
$$
that is,
$$
A_1 = aa' + cb', \quad B_1 = ba' + db', \quad C_1 = ac' + cd', \quad D_1 = bc' + dd',
$$
viz.\ $A_1$, $B_1$, $C_1$, $D_1$ are obtained from $A, B, C, D$ by the interchange of the accented and unaccented letters.

We have then, from the first expression for the Involutant,
\begin{align*}
I &= A \begin{vmatrix} 0 & b & b' \\ 0 & c & c' \\ 1 & d & d' \end{vmatrix}
- B \begin{vmatrix} 0 & c & c' \\ 1 & d & d' \\ 1 & a & a' \end{vmatrix}
+ C \begin{vmatrix} 1 & d & d' \\ 1 & a & a' \\ 0 & b & b' \end{vmatrix}
- D \begin{vmatrix} 1 & a & a' \\ 0 & b & b' \\ 0 & c & c' \end{vmatrix} \\
&= A(bc' - b'c) - B(ac' - a'c + cd' - c'd) + C(ab' - a'b + bd' - b'd) - D(bc' - b'c),
\end{align*}
or substituting for $A, B, C$ and $D$ their values, this is
\begin{multline*}
(aa' - dd' + bc' - b'c)(bc' - b'c) - (ab' + bd')(ac' - a'c + cd' - c'd) \\
+ (ca' + dc')(ab' - a'b + bd' - b'd) - (cb' + dd')(bc' - b'c);
\end{multline*}
and multiplying out and grouping together the terms in $bc$, $b'c'$, $bc'$ and $b'c$, this is found to be
$$
= -(a' - d')^2 bc + (a - d)(a' - d')(bc' + b'c) - (a - d)^2 b'c' + (bc' - b'c)^2,
$$
which is
$$
= -\{(a - d)b' - (a' - d')b\}\{(a - d)c - (a' - d')c'\} + (bc' - b'c)^2.
$$
Hence, writing
\begin{align*}
a &= bc' - b'c, \quad f = ad' - a'd, \\
b &= ca' - c'a, \quad g = bd' - b'd, \\
c &= ab' - a'b, \quad h = cd' - c'd,
\end{align*}
$$
I = -(c + g)(-b + h) + a^2 = bc - ch + bg - gh + a^2.
$$
To obtain the value of $I_1$, we must interchange the accented and unaccented letters, that is, change the signs of the several quantities $a, b, c, f, g, h$; but $I$, being a quadric function of the six quantities, is not altered by the change; that is, we have $I = I_1$.

\newpage
% ========================================================================
% PAPER 911: ON AN ALGEBRAICAL IDENTITY RELATING TO THE SIX COORDINATES OF A LINE
% ========================================================================

\section*{911. On an Algebraical Identity relating to the Six Coordinates of a Line.}
\addcontentsline{toc}{section}{911. On an Algebraical Identity relating to the Six Coordinates of a Line.}

\noindent\textit{[From the Messenger of Mathematics, vol.\ xx.\ (1891), pp.\ 138--140.]}

\vspace{0.5em}
\noindent\textbf{[p.\ 76]}

The following identity may be verified without difficulty; but it is interesting in regard to the analytical theory of the line.

The identity is
\begin{multline*}
0 = (bc' - b'c)^2(AF + BG + CH) \\
+ (bc' - b'c)(bC - cB)(Af + Bg + Ch + Fa + Gb + Hc) \\
- (bc' - b'c)(b'C - c'B)(Af' + Bg' + Ch' + Fa' + Gb' + Hc') \\
- (bC - cB)^2(af + bg + ch) \\
+ (bC - cB)(b'C - c'B)(af + a'f' + bg + b'g' + ch + c'h') \\
- (b'C - c'B)^2(a'f' + b'g' + c'h'),
\end{multline*}
where all the letters have arbitrary values.

It follows that, if
\begin{align*}
af + bg + ch &= 0, \\
a'f' + b'g' + c'h' &= 0, \\
af + a'f' + bg + b'g' + ch + c'h' &= 0, \\
AF + BG + CH &= 0, \\
Af + Bg + Ch + Fa + Gb + Hc &= 0, \\
Af' + Bg' + Ch' + Fa' + Gb' + Hc' &= 0,
\end{align*}

\vspace{0.5em}
\noindent\textbf{[p.\ 77]}

then either
$$
\begin{vmatrix} A & B & C \\ a & b & c \\ a' & b' & c' \end{vmatrix} = 0,
\quad \text{or else} \quad
\begin{vmatrix} B & C & F \\ b & c & f \\ b' & c' & f' \end{vmatrix} = 0.
$$

Supposing the first three of the six equations are satisfied, then $(a, b, c, f, g, h)$ and $(a', b', c', f', g', h')$ are the coordinates of two intersecting lines; and supposing that the last three equations are also satisfied, then $(A, B, C, F, G, H)$ will be the coordinates of a line meeting each of the intersecting lines. The third line is thus either in the plane of the intersecting lines, or else passes through their point of intersection; and in fact, the first of the two determinant equations is the condition in order that the line may be in the plane of the two intersecting lines, and the second determinant equation is the condition in order that it may pass through their point of intersection. Each equation is in fact one out of four equivalent forms, viz.\ we may have in the first equation $(B, C, F)$, $(C, A, G)$, $(A, B, H)$, or $(F, G, H)$; and in the second equation $(A, B, C)$, $(A, H, G)$, $(H, B, F)$, or $(G, F, C)$.

The analytical theory may be presented in a complete form by means of the formulae of my memoir ``On the six coordinates of a line,'' (1867), Camb.\ Phil.\ Trans., t.\ XI.\ pp.\ 290--323, [435]. Considering the two intersecting lines $(a, b, c, f, g, h)$ and $(a', b', c', f', g', h')$, the coordinates of the plane through these two lines (that is, the coefficients $\xi, \eta, \zeta, \omega$ of the equation $\xi x + \eta y + \zeta z + \omega w = 0$ of the plane) are there given (see p.\ 295*) in a fourfold form; and if we thence form the condition in order that the line $(A, B, C, F, G, H)$ may lie in this plane, we have
$$
\begin{pmatrix}
0 & C & -B & F \\
-C & 0 & A & G \\
B & -A & 0 & H \\
-F & -G & -H & 0
\end{pmatrix}
\begin{pmatrix}
af' + b'g + c'h & bf - b'f & cf - c'f & -(bc' - b'c) \\
ag' - a'g & a'f + bg' + c'h & cg' - c'g & -(ca' - c'a) \\
ah' - a'h & bh' - b'h & a'f + b'g + ch' & -(ab' - a'b) \\
gh' - g'h & hf - h'f & fg' - f'g & af + bg' + ch'
\end{pmatrix}
= 0;
$$
viz.\ the condition is expressible in any one of the 16 forms obtained by combining a line of the first matrix with a line of the second matrix; thus one form is
$$
0 \cdot (af' + b'g + c'h) + C(bf - b'f) - B(cf - c'f) - F(bc' - b'c) = 0,
$$
that is,
$$
\begin{vmatrix} B & C & F \\ b & c & f \\ b' & c' & f' \end{vmatrix} = 0,
$$
the foregoing first determinant equation, which thus belongs to the case where the line lies in the plane of the two intersecting lines.

\vspace{0.5em}
\noindent\textbf{[p.\ 78]}

Again we have (see p.\ 296*) an expression, in a fourfold form, for the coordinates of the point of intersection of the two intersecting lines; and thence for the condition, in order that the line $(A, B, C, F, G, H)$ may pass through this point, we have
$$
\begin{pmatrix}
0 & H & -G & A \\
-H & 0 & F & B \\
G & -F & 0 & C \\
-A & -B & -C & 0
\end{pmatrix}
\begin{pmatrix}
ag' - a'g & ah' - a'h & af + b'g + c'h & -(bc' - b'c) \\
bf - b'f & a'f + bg' + c'h & bh' - b'h & -(ca' - c'a) \\
cf - c'f & cg' - c'g & a'f + b'g + ch' & -(ab' - a'b) \\
gh' - g'h & hf - h'f & fg' - f'g & af + bg' + ch'
\end{pmatrix}
= 0;
$$
viz.\ the condition is expressible in any one of the 16 forms obtained by combining a line of the first matrix with a line of the second matrix; thus one form is
$$
A(bc' - b'c) + B(ca' - c'a) + C(ab' - a'b) + 0 \cdot (af + bg' + ch') = 0,
$$
that is,
$$
\begin{vmatrix} A & B & C \\ a & b & c \\ a' & b' & c' \end{vmatrix} = 0,
$$
the foregoing second determinant equation, which thus belongs to the case where the line passes through the point of intersection of the two intersecting lines.

I remark that the original identity may be written in the very compendious symbolical form
$$
\begin{vmatrix} A & a & a' \\ B & b & b' \\ C & c & c' \end{vmatrix}^2
\cdot
\begin{vmatrix} F & f & f' \\ G & g & g' \\ H & h & h' \end{vmatrix}^2
= 0,
$$
viz.\ here on the left-hand side the second factor denotes the three determinants $bc' - b'c$, $b'C - c'B$, $Bc - Cb$: the whole is a quadric function of these, the coefficients being
\begin{gather*}
AF + BG + CH, \quad af + bg + ch, \\
a'f' + b'g' + c'h', \quad af + bg' + ch' + fa' + gb' + hc',
\end{gather*}
$$
a'F + b'G + c'H + f'A + g'B + h'C, \quad Af + Bg + Ch + Fa + Gb + Hc,
$$
respectively; and the right-hand side is simply the product of two determinants.

\newpage
% ========================================================================
% PAPER 912: ON THE NOTION OF A PLANE CURVE OF A GIVEN ORDER
% ========================================================================

\section*{912. On the Notion of a Plane Curve of a Given Order.}
\addcontentsline{toc}{section}{912. On the Notion of a Plane Curve of a Given Order.}

\noindent\textit{[From the Messenger of Mathematics, vol.\ xx.\ (1891), pp.\ 148--150.]}

\vspace{0.5em}
\noindent\textbf{[p.\ 79]}

We have a complete geometrical notion of a curve of a given order, viz.\ a curve of the order $n$ is a curve which is met by any line whatever in $n$ points and no more; but starting with this definition, how do we know that there exists a curve of the order $n$? and, further, how do we know that it depends linearly on $\frac{1}{2}n(n + 3)$ parameters, or, what is the same thing, that there is one and only one curve which can be drawn through $\frac{1}{2}n(n + 3)$ given points?

The last-mentioned property does not by itself constitute a definition of a curve of the order $n$; thus $n = 2$, we cannot define a curve of the second order as a curve which is uniquely determined by the condition of passing through 5 given points; for a cubic passing through 4 given points is a curve uniquely determined by the condition in question; but we may differentiate between these two solutions by adding the further condition that, when 3 of the 5 points are in a line, the curve of the second order shall include as part of itself this line. And we are thus led to the definition:

\smallskip
\textit{A curve of the order $n$ is a curve which is uniquely determined by the condition of passing through $\frac{1}{2}n(n + 3)$ given points; and of being moreover such that, when $n + 1$ of these points lie on a line, it includes as part of itself this line.}
\smallskip

Starting from the foregoing definition, the first property is, I think, demonstrable, viz.\ the property that a curve of the order $n$ is met by any line whatever in $n$ points and no more. Thus $n = 2$: start with a line chosen at pleasure, and on it take 2 points which are regarded as indeterminate points: to fix the ideas, let one of these be regarded as depending on a parameter $\lambda$, and the other on a parameter $\mu$, (so that when $\lambda$ has a determinate value assigned to it, the first point becomes a determinate point, and similarly when $\mu$ has a determinate value assigned to it, the second point becomes a determinate point; and consequently, when $\lambda$ and $\mu$ have

\vspace{0.5em}
\noindent\textbf{[p.\ 80]}

determinate values, the 2 points are determinate points on the line). Take now any other 3 points not on the line; then, for the moment regarding the 2 points on the line as determinate points, we can through the 5 points draw a curve of the second order; this is the general curve of the second order through the 3 points; for it is a curve of the second order through the 3 points, and which, when the parameters $\lambda$ and $\mu$ are regarded as undetermined and arbitrary, or choosable at pleasure, might be made to pass through any other two points whatever. But by hypothesis, the curve meets the line in question, that is, any line, in two points: and the conic cannot meet the line in more than two points; for in like manner, starting with the given line, and upon it 3 points (which may be considered as depending on the parameters $\lambda$, $\mu$, and $\nu$ respectively), and taking any two points not on the line, we have through the 5 points a conic; and this conic, regarding the parameters $\lambda$, $\mu$, $\nu$ as undetermined and arbitrary, or choosable at pleasure, will be the general conic through the two points; for by a proper determination of the parameters, it might be made to pass through any other three points whatever. The reasoning is general, and we thus establish the theorem that a curve of the order $n$ is met by any line whatever in $n$ points and no more.

\vspace{1em}
\begin{center}
\rule{0.5\textwidth}{0.4pt}
\end{center}

\end{document}
